\chapter{Termodinamica}
\section{Variabili Termodinamiche}

I sistemi costituiti da un numero di elementi dell'ordine di $10^{23}$ sono irrisolvibili dal punto di vista pratico con sistemi di equazioni (es. $10^{23}$ equazioni). In questi sistemi (descritti microscopicamente) una minima variazione può compromettere le condizioni iniziali determinando un comportamento caotico. 

\medskip
Quindi dobbiamo sfruttare una \textbf{descrizione macroscopica} con pochi ($< 10$) \textbf{parametri globali}, il che la rende più semplice e predittiva, seppur meno precisa.

\begin{definizione}{Meccanica Statistica e Termodinamica}
La \textbf{meccanica statistica} stabilisce relazioni fra grandezze microscopiche e macroscopiche di tipo probabilistico.

\medskip
La \textbf{termodinamica} stabilisce relazioni fra grandezze macroscopiche su basi sperimentali tramite i principi termodinamici, aventi validità generale.
\end{definizione}

\subsection{Variabili Termodinamiche}
\begin{definizione}{Variabili termodinamiche}
Le variabili termodinamiche sono grandezze macroscopiche $A, B, C$ che descrivono globalmente un sistema tramite un'\textbf{equazione di stato} $F(A,B,C) = 0$.
\end{definizione}

\begin{itemize}
    \item \textbf{Massa totale} (sistema isolato) -- Meccaniche
    \item \textbf{Volume} (sistema deformabile) -- Meccaniche
    \item \textbf{Pressione} -- Nuova grandezza
    \item \textbf{Temperatura} -- Nuova grandezza
    \item \textbf{Quantità di calore} -- Nuova grandezza
\end{itemize}

Queste variabili sono \textbf{indipendenti dal tempo} e \textbf{definite in ogni punto}, cioè indipendenti dall'istante di partenza ma dipendenti dall'intervallo di tempo considerato.

\subsection{Pressione}
\begin{definizione}{Pressione}
La pressione è il rapporto fra il modulo della forza ortogonale a una superficie e l'area della superficie stessa:
\begin{equation}
p = \frac{||\vec{F}_{\perp}||}{S} = \frac{|\vec{F} \cdot \hat{n}|}{S}
\end{equation}
\end{definizione}

\medskip
\begin{itemize}
    \item Grandezza scalare
    \item Unità di misura nel SI: $1 \text{ Pa} = \frac{1 \text{ N}}{1 \text{ m}^2}$ (Pascal)
    \item $1 \text{ atm} = 101325 \text{ Pa} \simeq 101.3 \text{ kPa}$ (atmosfera)
    \item $1 \text{ bar} = 10^5 \text{ Pa}$
    \item $1 \text{ torr} = 133.3 \text{ Pa}$ (Torricelli)
\end{itemize}

\vspace{0.3cm}
\begin{figure}[ht]
\centering
\begin{tikzpicture}[scale=1]
    % Parallelogram representing the surface
    \fill[cyan!15, opacity=0.8] (0,0) -- (4,0) -- (5,1.5) -- (1,1.5) -- cycle;
    \draw[thick, mypen] (0,0) -- (4,0) -- (5,1.5) -- (1,1.5) -- cycle;
    \node[mypen, font=\bfseries] at (4.3, 0.4) {$\Delta S$};
    
    % Punto di applicazione
    \coordinate (P) at (2.5, 0.75);
    \filldraw[mypen] (P) circle (1.5pt);
    
    % Normal vector
    \draw[->, >=stealth, thick, mypen] (P) -- (2.5, 2.8) node[right] {$\hat{n}$};
    
    % Force vector
    \draw[->, >=stealth, thick, mygreen] (P) -- (3.8, 2.3) node[above right] {$\vec{F}$};
    
    % Proiezione normale di F
    \draw[dashed, mypen] (3.8, 2.3) -- (2.5, 2.3);
    \draw[->, >=stealth, thick, blue] (P) -- (2.5, 2.3) node[left, midway] {$F_\perp$};
\end{tikzpicture}
\caption{Rappresentazione vettoriale della forza $\vec{F}$ esercitata dal fluido su un elemento di superficie $\Delta S$; solo la componente normale $F_\perp$ definisce la pressione idrostatica.}
\end{figure}
\vspace{0.3cm}

\subsection{Stati di Aggregazione}
\begin{itemize}
    \item \textbf{Solido}: possiede forma e volume propri.
    \item \textbf{Fluido liquido}: possiede volume proprio, ma non forma propria.
    \item \textbf{Fluido aeriforme (gas)}: non possiede né volume né forma propri.
\end{itemize}

\begin{definizione}{Fluido Perfetto o Ideale}
Un fluido è detto ideale se presenta assenza di viscosità (attrito interno). In altre parole, considerando il fluido come un insieme di strati sovrapposti, quando due di questi scorrono l'uno sull'altro non si genera alcun attrito.
\end{definizione}
\begin{itemize}
    \item L'\textbf{acqua} è un fluido molto poco viscoso.
    \item L'\textbf{olio} è un fluido molto viscoso.
\end{itemize}

\section{Pressione in un Fluido Ideale}
In un fluido perfetto all'equilibrio, preso un punto P e una superficie $\Delta S$ intorno a P, nel limite per $\Delta S \to 0$ la pressione $F/\Delta S$ è la stessa in ogni direzione (Principio di Pascal).

\begin{figure}[ht]
    \centering
    \begin{tikzpicture}[scale=1.1]
        % Faccia frontale (triangolo rettangolo)
        \fill[cyan!15, opacity=0.8] (0,0) -- (3,0) -- (0,2) -- cycle;
        \draw[thick, mypen] (0,0) -- (3,0) -- (0,2) -- cycle;
        
        % Spigoli posteriori (prospettiva)
        \draw[dashed, mypen] (0,0) -- (1,0.6);
        \draw[dashed, mypen] (3,0) -- (4,0.6);
        \draw[dashed, mypen] (0,2) -- (1,2.6);
        \draw[dashed, mypen] (1,0.6) -- (4,0.6) -- (1,2.6) -- cycle;
        
        % Spigoli visibili superiori
        \draw[thick, mypen] (3,0) -- (4,0.6) -- (1,2.6) -- (0,2);
        
        % Etichetta punto P interno
        \filldraw[mypen] (1.5,0.8) circle (1.5pt) node[below right] {P};
        
        % Vettori pressione sulle facce
        \draw[->, >=stealth, thick, mygreen] (-0.8, 1) -- (0, 1) node[midway, above] {$\vec{F}_x$};
        \draw[->, >=stealth, thick, mygreen] (1.5, -0.8) -- (1.5, 0) node[midway, left] {$\vec{F}_y$};
        \draw[->, >=stealth, thick, notered] (2.8, 1.8) -- (1.8, 1.1) node[midway, above right] {$\vec{F}_n$};
    \end{tikzpicture}
    \caption{Prisma infinitesimo in equilibrio in un fluido ideale: la dimostrazione del Principio di Pascal mostra l'isotropia della pressione nel limite per il volume che tende a zero.}
\end{figure}

\medskip
Se la regione A esercita $\vec{F}$ su $\Delta S$, allora per il terzo principio della dinamica la regione B esercita $-\vec{F}$ su $\Delta S$.
\begin{itemize}
    \item Se il fluido è \textbf{viscoso}, allora $\vec{F} \not\parallel \hat{n}$.
    \item Se il fluido è \textbf{ideale}, allora $\vec{F} \parallel \hat{n}$.
\end{itemize}

All'equilibrio la somma delle forze applicate sul prisma su ogni sua superficie è nulla: $\sum \vec{F}_s = 0$.
Sia $\vec{P} = -mg \hat{z}$ il peso del fluido:
\begin{equation}
\vec{F}_S + \vec{F}_1 + \vec{F}_2 + \vec{F}_3 + \vec{F}_4 + \vec{P} = 0
\end{equation}

Scomponendo lungo gli assi:
\begin{equation*}
\begin{cases}
    \hat{x}: & F_4 - F_3 = 0 \\
    \hat{y}: & p_y b^2 \sin\theta - p b^2 \sin\theta = 0 \\
    \hat{z}: & p_z b^2 \cos\theta - p b^2 \cos\theta - mg = 0
\end{cases}
\implies
\begin{cases}
    p_y = p \\
    p_z = p + \frac{mg}{b^2 \cos\theta}
\end{cases}
\end{equation*}

Poiché $mg = \rho V_{\text{prisma}} g = \rho \left( \frac{b^2 \sin\theta \cos\theta \cdot b}{2} \right) g = \rho \frac{b^3 \sin\theta \cos\theta}{2} g$, sostituendo in $p_z$:
\begin{equation}
p_z = p + \frac{\rho b^3 \sin\theta \cos\theta g}{2 b^2 \cos\theta} = p + \frac{1}{2} \rho b \sin\theta g
\end{equation}

Nel limite per $\Delta S \to 0$, che equivale a far tendere $b \to 0$:
\begin{equation}
\lim_{b \to 0} \left( p + \frac{1}{2} \rho b \sin\theta g \right) = p
\end{equation}

Di conseguenza, ruotando il prisma di $90^\circ$ si dimostra analogamente per l'altro asse che:
\begin{equation}
\boxed{p = p_y = p_z = p_x}
\end{equation}

\vspace{0.3cm}
\section{Legge di Stevino}
\begin{definizione}{Legge di Stevino}
In un fluido all'equilibrio, la pressione esercitata su uno strato del fluido è determinata soltanto dal peso degli strati sovrastanti (che si trovano a una profondità $z$ minore).
\end{definizione}

\medskip
All'equilibrio $\sum \vec{F}_z = 0$. Sia $\vec{R}$ la forza di reazione degli strati a $z' > z$ e $\vec{P}$ la forza peso degli strati a $z' \le z$. Si ottiene:
\begin{equation}
\vec{P} + \vec{R} = 0 \implies \vec{R} = -m\vec{g} \implies R = mg = \rho V g = \rho S z g = p S
\end{equation}

Da cui la pressione:
\begin{equation}
p = \frac{R}{S} = \frac{\rho g S z}{S} = \rho g z
\end{equation}

Se si considera anche la pressione atmosferica o esterna esercitata al livello $z=0$, allora si ottiene la formula generale:
\begin{equation}
\boxed{p(z) = p(0) + \rho g z}
\end{equation}

\vspace{0.3cm}
\section{Legge di Archimede}
\begin{definizione}{Spinta di Archimede}
Un corpo immerso in un fluido riceve una spinta dal basso verso l'alto pari al peso del volume di fluido spostato dal corpo.
\end{definizione}

\medskip
La forza risultante lungo l'asse $z$ si ricava come differenza tra la pressione sulla faccia inferiore e quella sulla faccia superiore:
\begin{align*}
F_z &= F_{\text{up}} - F_{\text{down}} = p(z_2)S - p(z_1)S - mg \\
&= (p(0) + \rho_{\text{fluido}} g z_2)S - (p(0) + \rho_{\text{fluido}} g z_1)S - mg \\
&= \rho_{\text{fluido}} g S (z_2 - z_1) - mg \\
&= \rho_{\text{fluido}} g S (z_2 - z_1) - \rho_{\text{corpo}} g S (z_2 - z_1) \\
&= (\rho_{\text{fluido}} - \rho_{\text{corpo}}) g S (z_2 - z_1) = (\rho_{\text{fluido}} - \rho_{\text{corpo}}) g V_{\text{corpo}}
\end{align*}

Dunque la forza risultante (tenendo conto anche del peso del corpo) è:
\begin{equation}
\boxed{F_z = (\rho_{\text{fluido}} - \rho_{\text{corpo}}) g V_{\text{corpo}}}
\end{equation}

\begin{itemize}
    \item Se $\rho_{\text{fluido}} < \rho_{\text{corpo}} \implies \vec{F}_z \parallel -\hat{z}$, cioè il corpo va a fondo.
    \item Se $\rho_{\text{fluido}} > \rho_{\text{corpo}} \implies \vec{F}_z \parallel +\hat{z}$, cioè il corpo galleggia.
\end{itemize}

Nel caso di galleggiamento, il corpo risale finché non raggiunge una situazione di equilibrio in cui la forza totale è nulla ($F_z = 0$). L'equilibrio si ha quando:
\begin{equation}
\rho_{\text{fluido}} g V_{\text{immerso}} = mg
\end{equation}
dove $V_{\text{immerso}}$ è il volume della parte del corpo effettivamente immersa nel fluido. Pertanto, uguagliando le masse e considerando l'altezza d'immersione:
\begin{equation}
\rho_{\text{fluido}} S (z_2 - \bar{z}) = \rho_{\text{corpo}} S (z_2 - z_1) \implies \boxed{(z_2 - \bar{z}) = (z_2 - z_1) \frac{\rho_{\text{corpo}}}{\rho_{\text{fluido}}}}
\end{equation}
che rappresenta la profondità di immersione.

\section{Misura della Pressione}
La pressione si misura con i \textbf{manometri} (ad es. per gli pneumatici) e i \textbf{barometri} (ad es. per l'atmosfera).
La pressione atmosferica fu misurata per la prima volta da Evangelista Torricelli nel 1643, il quale inventò il barometro.

\begin{figure}[ht]
    \centering
    \begin{tikzpicture}
        % Barometro di Torricelli
        \draw[thick, mypen] (0,0) -- (3,0) -- (3,1) -- (0,1) -- cycle; % vaschetta
        \draw[thick, mypen, fill=blue!30] (0.1,0.1) rectangle (2.9,0.8); % liquido in vaschetta
        \draw[thick, mypen] (1.3,0.5) -- (1.3,4) -- (1.7,4) -- (1.7,0.5); % tubo
        \draw[thick, mypen, fill=blue!30] (1.3,0.5) rectangle (1.7,3); % liquido in tubo
        \node at (2.5,0.5) {A \quad B};
        \node at (1.5,3.5) {vuoto};
        \draw[->, >=stealth, mypen] (3.5, 2) -- (1.8, 3) node[right, xshift=1cm] {aria};
    \end{tikzpicture}
    \caption{Barometro di Torricelli}
\end{figure}

\medskip
L'esperimento utilizza il mercurio e si basa sui seguenti principi:
\begin{itemize}
    \item La densità del mercurio è $\rho_{Hg} = 13.7 \cdot 10^3 \text{ kg/m}^3$ ($\simeq 14 \rho_{H_2O}$), essendo un metallo liquido molto pesante.
    \item L'altezza della colonna di mercurio all'equilibrio è $h_3 = 76 \text{ cm}$.
    \item $p_A = p_B$: se le pressioni fossero diverse ci sarebbero forze non nulle che farebbero entrare o uscire mercurio dal capillare, ma essendo il sistema all'equilibrio le pressioni si equivalgono.
    \item $p_B = \rho_{Hg} g h_1 + \rho_{Hg} g h_2 + p_{\text{atm}}$
    \item $p_A = \rho_{Hg} g h_1 + \rho_{Hg} g h_2 + \rho_{Hg} g h_3$
\end{itemize}

Uguagliando le due pressioni si ottiene:
\begin{align*}
p_{\text{atm}} &= \rho_{Hg} g h_3 = 13.7 \cdot 10^3 \cdot 9.81 \cdot 0.76 \text{ Pa} \simeq 101 \text{ kPa} \\
&\implies \boxed{p_{\text{atm}} = \rho_{Hg} g h_3 \simeq 101 \text{ kPa} \simeq 760 \text{ Torr}}
\end{align*}

\begin{osservazione}{Il peso dell'aria}
L'aria, per quanto impercettibile, ha un peso e si può "maneggiare", cioè si può mettere e togliere: lo spazio vuoto che si viene a creare nel capillare quando scende il mercurio ne è la dimostrazione.
\end{osservazione}


\section{Temperatura}
\begin{definizione}{Temperatura}
La temperatura è la misura della variazione di volume (o altra proprietà termometrica) di un corpo posto a contatto con una sorgente di calore.
\end{definizione}

\medskip
La misura si effettua con i \textbf{termometri}. Per tararli, si scelgono dei punti fissi $P_1$ e $P_2$ a cui si attribuiscono determinati valori di temperatura; questi punti devono essere \textbf{facilmente riproducibili} e \textbf{controllabili}.

A pressione $p_{\text{atm}}$ costante, i punti di riferimento scelti sono:
\begin{itemize}
    \item $P_1$: punto di fusione dell'acqua (equilibrio tra acqua e ghiaccio fondente).
    \item $P_2$: punto di ebollizione dell'acqua (equilibrio tra acqua e vapore).
\end{itemize}

\subsection{Scala Celsius}
Vengono assegnati i valori $P_1 \to 0^\circ\text{C}$ e $P_2 \to 100^\circ\text{C}$.

\subsection{Scala Fahrenheit}
Vengono assegnati i valori $P_1 \to 32^\circ\text{F}$ e $P_2 \to 212^\circ\text{F}$.

\begin{figure}[ht]
    \centering
    \begin{tikzpicture}
        \draw[->, >=stealth, mypen] (0,0) -- (5,0) node[right] {$^\circ\text{C}$};
        \draw[->, >=stealth, mypen] (0,0) -- (0,3) node[above] {$^\circ\text{F}$};
        \draw[mypen] (0,0.5) node[left] {32} -- (0.1,0.5);
        \draw[mypen] (0,2.5) node[left] {212} -- (0.1,2.5);
        \draw[mypen] (3,0) node[below] {100} -- (3,0.1);
        \draw[mypen] (0,0) node[below right] {0};
        \draw[thick, blue] (-1,-0.1) -- (4,3.1);
    \end{tikzpicture}
    \caption{Relazione tra le scale Celsius e Fahrenheit}
\end{figure}

\medskip
La formula per la conversione da $^\circ\text{C}$ a $^\circ\text{F}$ è la seguente:
\begin{equation}
    \boxed{T(^\circ\text{F}) = 32 + \frac{9}{5} T(^\circ\text{C})}
\end{equation}

\section{Sostanze Termometriche e Termometri}
Come sostanza di cui misurare la variazione di volume si usano generalmente gas o liquidi, per la loro elevata capacità di dilatarsi termicamente.

\subsection{Termometro a Mercurio}
Si segna l'altezza della colonna di mercurio quando è inserita nell'acqua bollente ($100^\circ\text{C}$) e nel ghiaccio fondente ($0^\circ\text{C}$). Successivamente si divide in 100 intervalli identici per creare la \textbf{scala centigrada}.

Affinché il termometro sia reattivo, il volume del capillare deve essere molto minore di quello del bulbo ($V_{\text{capillare}} \ll V_{\text{bulbo}}$) e il bulbo non può essere troppo grande, altrimenti il liquido impiega troppo tempo per dilatarsi uniformemente.

Nella scelta del fluido termometrico si deve tenere conto anche delle sue temperature di congelamento ed ebollizione:
\begin{itemize}
    \item \textbf{Mercurio (Hg)}: congelamento a $-39^\circ\text{C}$ (si solidificherebbe con un clima artico) ed ebollizione a $357^\circ\text{C}$.
    \item \textbf{Alcol}: congelamento a $-70^\circ\text{C}$ ed ebollizione a $78.5^\circ\text{C}$ (si vaporizzerebbe prima dell'acqua).
\end{itemize}

\vspace{0.3cm}
\section{Principio Zero della Termodinamica}
\begin{definizione}{Principio Zero}
Due corpi che sono separatamente in equilibrio termico con un terzo corpo, sono in equilibrio termico anche tra di loro.
\end{definizione}

\begin{figure}[ht]
    \centering
    \begin{tikzpicture}[scale=0.95]
        % Parete adiabatica esterna e di separazione tra A e B
        \fill[pattern=north east lines, mypen!30] (0,0) rectangle (6,3);
        \fill[white] (0.2,0.2) rectangle (5.8,2.8);
        \fill[pattern=north east lines, mypen!30] (2.9,1.4) rectangle (3.1,2.8);
        
        % Sistema A e B
        \draw[thick, mypen, fill=cyan!10] (0.5,1.5) rectangle (2.8,2.6) node[midway, font=\bfseries] {Sistema A ($T_A$)};
        \draw[thick, mypen, fill=orange!10] (3.2,1.5) rectangle (5.5,2.6) node[midway, font=\bfseries] {Sistema B ($T_B$)};
        
        % Sistema C (Termometro/riferimento in contatto diatermico con A e B)
        \draw[thick, mypen, fill=green!10] (0.5,0.4) rectangle (5.5,1.3) node[midway, font=\bfseries] {Sistema C ($T_C$) - Parete Diatermica};
        
        % Frecce di equilibrio termico
        \draw[<->, >=stealth, thick, notered] (1.65, 1.4) -- (1.65, 1.6) node[midway, left] {Eq.};
        \draw[<->, >=stealth, thick, notered] (4.35, 1.4) -- (4.35, 1.6) node[midway, right] {Eq.};
    \end{tikzpicture}
    \caption{Illustrazione del Principio Zero: i sistemi A e B, separati da una parete adiabatica, sono singolarmente in equilibrio termico con il terzo sistema C (grazie alla parete diatermica). Ne consegue che A e B sono in equilibrio termico tra loro ($T_A = T_B = T_C$).}
\end{figure}

All'inizio si ha $T_1 \neq T_2$; se li poniamo a contatto termico fra loro (isolandoli dall'esterno), dopo un certo tempo $\Delta t$ essi raggiungono la stessa temperatura finale $T = T_1' = T_2'$, trovandosi in uno stato di \textbf{equilibrio termico}.
Grazie alla proprietà transitiva, se $T_A = T_B$ e $T_B = T_C$, allora $T_A = T_C$.

La temperatura è una grandezza fondamentale, la cui unità di misura nel SI è il Kelvin (K).

\vspace{0.3cm}
\section{Termometro a Gas}
Ci sono termometri a gas di due tipi: a pressione o a volume costante. Nonostante la maggiore precisione, sono anche più ingombranti dei termometri a liquido.

\subsection{Termometro a Pressione Costante}
Un rubinetto collega due capillari fra loro e con un serbatoio (isolato). Se lasciamo il rubinetto aperto e scaldiamo il gas, questo si dilata e spinge il mercurio dall'altra parte.
Quindi facciamo defluire il mercurio nel serbatoio in modo da azzerare il dislivello ($h_1 - h_2 = 0$).

All'equilibrio, la pressione nel punto A può essere calcolata da entrambe le parti:
\begin{align*}
p_A(\text{Dx}) &= p_{\text{gas}} + \rho_{Hg} g h_2 \\
p_A(\text{Sx}) &= p_{\text{atm}} + \rho_{Hg} g h_1
\end{align*}

Dunque si ha la pressione $p_{\text{gas}} = p_{\text{atm}}$.

Misurando il volume del gas si ottiene:
\begin{equation}
    \boxed{V(T_C) = V_0 + \alpha V_0 T_C}
\end{equation}
dove $V_0$ è il volume a $T_C=0^\circ\text{C}$ e $p=1\text{ atm}$.
Il coefficiente $\alpha$ calcolato sperimentalmente vale:
\begin{equation}
    \alpha = \frac{V_{100} - V_0}{100 V_0} (^\circ\text{C})^{-1} = \boxed{ \left( 273.15 \pm 0.01 \, ^\circ\text{C} \right)^{-1} \simeq \frac{1}{273.15} \, (^\circ\text{C})^{-1} }
\end{equation}
Questo valore è comune a gran parte dei gas a pressione atmosferica e temperatura molto più elevata di quella di liquefazione.

\subsection{Termometro a Volume Costante}
Alzando e abbassando il serbatoio di mercurio a seconda dei casi variamo la pressione in funzione della temperatura, mantenendo il volume costante:
\begin{equation}
    \boxed{p(T_C) = p_0 + \alpha p_0 T_C}
\end{equation}
dove $p_0$ è la pressione a $T_C=0^\circ\text{C}$. Il coefficiente $\alpha$ risulta identico, comprensivo dell'incertezza sperimentale: $\alpha = \frac{p_{100} - p_0}{100 p_0} (^\circ\text{C})^{-1} = \boxed{ \left( 273.15 \pm 0.01 \, ^\circ\text{C} \right)^{-1} }$.

\vspace{0.3cm}
\section{Scala Kelvin e Gas Perfetti}
\begin{definizione}{Gas Perfetto}
Le due formule trovate con i termometri a gas a pressione e volume costante sono chiamate \textbf{Leggi di Gay-Lussac}. Un gas che segue queste due leggi per qualsiasi $p, T$ è detto \textbf{gas perfetto} o ideale.
\end{definizione}

Il gas perfetto in natura non esiste, ma molti dei gas nobili hanno un comportamento simile in ampi intervalli di $p$ e $T$.

\medskip
Se calcoliamo i valori per cui $V(T_C) = 0$ e $p(T_C) = 0$, otteniamo $T_C = -1/\alpha \simeq -273.15^\circ\text{C}$.
Esiste quindi una temperatura al di sotto della quale la materia si comporta in modo differente. Si definisce pertanto una \textbf{scala assoluta} detta \textbf{scala Kelvin} che ha come valore $0\text{ K} = -1/\alpha$, detto \textbf{zero assoluto}.

\begin{equation}
\boxed{T(\text{K}) = T(^\circ\text{C}) + 273.15}
\end{equation}
La conversione da $^\circ\text{C}$ a K mantiene lo stesso intervallo: $1^\circ\text{C} = 1\text{ K}$.

Scrivendo le leggi di Gay-Lussac rispetto alla scala Kelvin si ottiene:
\begin{align}
    V(T) &= \alpha V_0 T \quad \text{a } p=\text{costante} \\
    p(T) &= \alpha p_0 T \quad \text{a } V=\text{costante}
\end{align}

Per un gas perfetto a $T=0\text{ K} \implies p=0, V=0$.
Con la meccanica statistica si può dimostrare che l'energia cinetica media $\langle K \rangle \propto T$. Ciò significa che per $T \to 0$ il sistema non può scambiare energia cedendola a un altro sistema, al massimo può solo riceverla, poiché ha raggiunto l'energia minima possibile.
Con "scala assoluta" inoltre si intende che essa è indipendente dalla sostanza termometrica (scala $\Theta$, legata al Secondo Principio della Termodinamica).

\vspace{0.3cm}
\subsection{Derivazione dell'Equazione di Stato, Condizioni STP e Legge di Dalton}
Combinando le due leggi di Gay-Lussac con la legge di Boyle-Mariotte ($p V = \text{costante}$ a $T = \text{costante}$), si ricava l'equazione di stato generale che collega le variabili termodinamiche di una massa gassosa. Consideriamo di passare da uno stato iniziale $A(p_0, V_0, T_0)$ a uno stato finale $B(p, V, T)$ attraverso uno stato intermedio ausiliario $A^*(p_0, V^*, T)$ mediante una trasformazione isobara ($A \to A^*$) e una successiva isoterma ($A^* \to B$):
\begin{align}
    \text{Isobara } A \to A^* &: \quad \frac{V^*}{V_0} = \frac{T}{T_0} \implies V^* = V_0 \frac{T}{T_0} \\
    \text{Isoterma } A^* \to B &: \quad p V = p_0 V^* \implies p V = p_0 V_0 \frac{T}{T_0} \implies \frac{p V}{T} = \frac{p_0 V_0}{T_0} = \text{costante}
\end{align}
Per una mole di gas ($n = 1\text{ mol}$), tale costante universale è indicata con $R$ (\textbf{Costante Universale dei Gas}). Generalizzando a un numero di moli $n = m / M_{\text{molare}}$, si ottiene l'\textbf{Equazione di Stato dei Gas Perfetti}:
\begin{equation}
    \boxed{ p V = n R T = \frac{N}{N_A} R T = N k_B T }
\end{equation}
dove $N$ è il numero totale di molecole, $N_A = 6.022 \cdot 10^{23}\text{ mol}^{-1}$ è il \textbf{Numero di Avogadro} e $k_B = R / N_A = 1.381 \cdot 10^{-23}\text{ J/K}$ è la \textbf{Costante di Boltzmann}.

\begin{definizione}{Condizioni STP e Legge di Avogadro}
Le \textbf{Condizioni Standard di Temperatura e Pressione (STP)} corrispondono a $T_0 = 273.15\text{ K}$ ($0^\circ\text{C}$) e $p_0 = 101.325\text{ kPa} = 1\text{ atm}$. 
Per la \textbf{Legge di Avogadro}, volumi uguali di gas diversi, alle medesime condizioni di temperatura e pressione, contengono lo stesso numero di molecole. Nelle condizioni STP, una mole ($n=1\text{ mol}$) di qualsiasi gas perfetto occupa un volume molare standard pari a:
\begin{equation}
    V_{\text{mol}} = \frac{R T_0}{p_0} = 22.414\text{ L} = 22.414 \cdot 10^{-3}\text{ m}^3
\end{equation}
I valori numerici della costante universale dei gas $R$ nelle diverse unità di misura sono pertanto:
\begin{equation}
    R = 8.314 \, \frac{\text{J}}{\text{mol} \cdot \text{K}} = 0.08205 \, \frac{\text{L} \cdot \text{atm}}{\text{mol} \cdot \text{K}} = 1.987 \, \frac{\text{cal}}{\text{mol} \cdot \text{K}}
\end{equation}
\end{definizione}

\begin{definizione}{Legge di Dalton delle Pressioni Parziali}
In un recipiente di volume $V$ contenente una miscela di k gas chimicamente inerti a temperatura $T$, la pressione totale esercitata dalla miscela sulle pareti è uguale alla somma delle \textbf{pressioni parziali} che ciascun gas eserciterebbe se occupasse da solo l'intero volume:
\begin{equation}
    \boxed{ p_{\text{tot}} = \sum_{i=1}^k p_i = \sum_{i=1}^k \frac{n_i R T}{V} = \frac{R T}{V} \sum_{i=1}^k n_i = \frac{n_{\text{tot}} R T}{V} }
\end{equation}
La pressione parziale dell'i-esimo componente è direttamente proporzionale alla sua frazione molare $x_i = n_i / n_{\text{tot}}$, ovvero $p_i = x_i \, p_{\text{tot}}$.
\end{definizione}

\vspace{0.3cm}
\section{Punto Triplo}
\begin{definizione}{Punto Triplo}
Il punto triplo è la combinazione di pressione e temperatura in cui in una sostanza coesistono i tre stati di aggregazione della materia in equilibrio termodinamico.
\end{definizione}

\medskip
Per l'acqua ($\text{H}_2\text{O}$) questo si verifica a $T \simeq 0.01^\circ\text{C}$ (temperatura alla quale \textbf{deve esistere il ghiaccio fondente}) e $p \simeq 4.58 \text{ Torr}$ (pressione alla quale \textbf{deve esistere l'acqua bollente}): la simultaneità di tali coordinate fisiche garantisce la coesistenza in equilibrio dinamico di ghiaccio, acqua liquida e vapore.
L'unità di misura K (Kelvin) della scala termodinamica assoluta si definisce come la frazione $1/273.16$ della temperatura termodinamica del punto triplo dell'acqua:
\begin{equation}
    \boxed{ 1\text{ K} = \frac{T_{\text{punto triplo}}}{273.16} }
\end{equation}
il fatto che sia 16 deriva dal fatto che c'è una miscela differente di isotopi per i quali il punto di congelamento differisce lievemente.

\begin{osservazione}{Significato quantistico di $0\text{ K}$}
Considerando i livelli quantizzati di energia $E_0, E_1, E_2, \dots$ di un sistema, lo \textbf{stato fondamentale} (di minima energia) corrisponde a $E_0$.
Per un sistema di $N$ particelle si ha quindi $E_{\text{tot}} = N \cdot E_0 = E_{0,\text{minima}}$.
Appena una particella passa a uno stato eccitato ($E_i > E_0$) allora $E_{\text{tot}} > E_0$.
Se il sistema è isolato in un bagno termico a temperatura $T=0\text{ K}$, non vi è alcuno scambio termico verso l'esterno che possa abbassare ulteriormente l'energia, né dall'esterno che la innalzi. Si conclude quindi che $0\text{ K}$ è la temperatura a cui ogni sistema si trova nel suo stato fondamentale.
\end{osservazione}

\vspace{0.3cm}
\section{Dilatazione Termica}
\begin{definizione}{Dilatazione termica}
La dilatazione termica è la variazione di ogni dimensione (lineare, superficiale, volumica) di un corpo al variare della temperatura ($T \uparrow \implies V \uparrow$).
\end{definizione}

\medskip
Se $V(0^\circ\text{C}) = V_0$, la variazione può essere descritta come $V(T_C) = V_0 F(T_C)$, dove la funzione $F(0) = 1$ è adimensionale.
Sviluppando in serie di Taylor nell'intorno dello zero si ottiene:
\begin{equation}
F(T_C) = F(0) + \left(\frac{\partial F}{\partial T_C}\right)_0 T_C + \frac{1}{2} \left(\frac{\partial^2 F}{\partial T_C^2}\right)_0 T_C^2 + o(T_C^3) \simeq 1 + \beta_0 T_C
\end{equation}
in quanto sperimentalmente si osserva che la relazione è lineare. Di conseguenza:
\begin{align}
\beta_0 &= \left(\frac{\partial F}{\partial T_C}\right)_0 = \frac{1}{V_0} \left(\frac{\partial V(T_C)}{\partial T_C}\right)_0 
\end{align}
rappresenta il \textbf{coefficiente di dilatazione termica} a $T_C = 0^\circ\text{C}$ e pressione costante.

\subsection{Dilatazione Termica Lineare ($\beta_0 = \lambda$)}
% Disegno di una sbarra di lunghezza l
\begin{figure}[ht]
\centering
\begin{tikzpicture}[scale=1]
\fill[pattern=north east lines, mypen!20] (0,0) rectangle (4,0.3);
\draw[thick, mypen] (0,0) rectangle (4,0.3);
\draw[<->, >=stealth, thick, notered] (0,0.6) -- (4,0.6) node[midway, above, font=\bfseries] {$\ell(T)$};
\draw[dashed, mypen] (0,0.3) -- (0,0.8);
\draw[dashed, mypen] (4,0.3) -- (4,0.8);
\end{tikzpicture}
\caption{Dilatazione termica lineare di una sbarra metallica all'aumentare della temperatura.}
\end{figure}

Il coefficiente lineare $\lambda$ è definito come $\lambda = \frac{1}{\ell_0} \left. \left( \frac{\partial \ell}{\partial T_C} \right) \right|_{T_C = 0^\circ\text{C}}$.
\begin{equation}
\frac{\partial \ell}{\ell} = \lambda \partial T_C \implies \int_{\ell_0}^{\ell(T_C)} \frac{\partial \ell}{\ell} = \ln \left( \frac{\ell(T_C)}{\ell_0} \right) = \int_0^{T_C} \lambda \partial T_C' = \lambda T_C
\end{equation}

Da cui:
\begin{equation}
\ell(T_C) = \ell_0 e^{\lambda T_C} \stackrel{\lambda T_C \ll 1}{\simeq} \ell_0 (1 + \lambda T_C)
\end{equation}
Nel campo di esistenza dei solidi questa approssimazione è accurata. 
Inoltre $\Delta \ell = \ell_0 \lambda T_C \implies \frac{\Delta \ell}{\ell_0} \ll 1 \implies \Delta \ell \ll \ell_0$ poiché $T_C \ll \lambda^{-1} \ll T_{\text{fusione}}$.

\subsection{Dilatazione Termica Superficiale ($\beta_0 = 2\lambda$)}
% Disegno di un rettangolo a x b
\begin{figure}[ht]
\centering
\begin{tikzpicture}[scale=1.1]
\fill[cyan!15, opacity=0.8] (0,0) rectangle (2.5,1.8);
\draw[thick, mypen] (0,0) rectangle (2.5,1.8);
\draw[<->, >=stealth, mypen] (0,-0.2) -- (2.5,-0.2) node[midway, below] {$a$};
\draw[<->, >=stealth, mypen] (-0.2,0) -- (-0.2,1.8) node[midway, left] {$b$};
\node[mypen, font=\bfseries] at (1.25, 0.9) {$S = a b$};
\end{tikzpicture}
\caption{Dilatazione termica superficiale di una lamina rettangolare: la variazione di area è proporzionale al coefficiente $\beta_0 = 2\lambda$.}
\end{figure}

Sviluppando in serie e trascurando i termini di ordine superiore:
\begin{align*}
S(T_C) &= a(T_C) \cdot b(T_C) = a_0 b_0 (1 + \lambda T_C)^2 \\
&= S_0 (1 + \lambda T_C)^2 = S_0 (1 + \lambda^2 T_C^2 + 2\lambda T_C) \\
&\simeq S_0 (1 + 2\lambda T_C)
\end{align*}
da cui si evince che il coefficiente di dilatazione superficiale è $2\lambda$.

\subsection{Dilatazione Termica Volumica ($\beta_0 = 3\lambda$)}
% Disegno di un parallelepipedo a x b x c
\begin{figure}[ht]
\centering
\begin{tikzpicture}[scale=1]
% Parallepipedo in prospettiva
\fill[orange!15, opacity=0.8] (0,0) -- (2.5,0) -- (2.5,1.5) -- (0,1.5) -- cycle;
\fill[orange!25, opacity=0.8] (0,1.5) -- (0.8,2.1) -- (3.3,2.1) -- (2.5,1.5) -- cycle;
\fill[orange!35, opacity=0.8] (2.5,0) -- (3.3,0.6) -- (3.3,2.1) -- (2.5,1.5) -- cycle;
\draw[thick, mypen] (0,0) -- (2.5,0) -- (2.5,1.5) -- (0,1.5) -- cycle;
\draw[thick, mypen] (0,1.5) -- (0.8,2.1) -- (3.3,2.1) -- (2.5,1.5);
\draw[thick, mypen] (2.5,0) -- (3.3,0.6) -- (3.3,2.1);
\draw[dashed, mypen] (0,0) -- (0.8,0.6) -- (3.3,0.6);
\draw[dashed, mypen] (0.8,0.6) -- (0.8,2.1);
\node[below] at (1.25, 0) {$a$};
\node[right] at (3.3, 1.35) {$c$};
\node[above left] at (0.4, 1.8) {$b$};
\node[mypen, font=\bfseries] at (4.5, 1) {$V = a b c$};
\end{tikzpicture}
\caption{Dilatazione termica volumica di un solido geometrico: la dilatazione tridimensionale è regolata da $\beta_0 = 3\lambda$.}
\end{figure}

Sviluppando analogamente si ottiene:
\begin{align*}
V(T_C) &= a(T_C) \cdot b(T_C) \cdot c(T_C) = a_0 b_0 c_0 (1 + \lambda T_C)^3 \\
&= V_0 (1 + \lambda T_C)^3 = V_0 (1 + 3\lambda^2 T_C + 3\lambda T_C^2 + \lambda^3 T_C^3) \\
&\simeq V_0 (1 + 3\lambda T_C)
\end{align*}
da cui si evince che il coefficiente di dilatazione volumica è $3\lambda$.

\vspace{0.3cm}
\subsection{Interpretazione Quantitativa a Livello Molecolare}
% Grafico curva dell'energia potenziale molecolare
\begin{figure}[ht]
\centering
\begin{tikzpicture}[scale=1]
\draw[->, >=stealth, mypen] (-0.5, 0) -- (4, 0) node[right] {$r$};
\draw[->, >=stealth, mypen] (0, -3) -- (0, 2) node[above] {$U(r)$};
\draw[thick, blue] plot [smooth] coordinates {(0.2, 1.5) (0.3, 0) (0.5, -2.5) (1, -1.5) (2, -0.5) (3.5, -0.1)};
\draw[dashed, mypen] (0.5, -2.5) -- (0, -2.5) node[left] {$E_{min}$};
\draw[dashed, mypen] (0.5, -2.5) -- (0.5, 0) node[above] {$r_0$};
\draw[dashed, mypen] (0.15, -1) -- (1.5, -1);
\draw[dashed, mypen] (0.2, -1.8) -- (1, -1.8);
\node at (-0.3, -1) {$E_2$};
\node at (-0.3, -1.8) {$E_1$};
\node at (0.8, -1.6) {$P_1$};
\node at (1.2, -0.8) {$P_2$};
\node at (0.5, -2.7) {$P_0$};
\end{tikzpicture}
\caption{Curva dell'energia potenziale molecolare}
\end{figure}
Sia $U_{\text{eff}}(r)$ l'energia potenziale nel caso generale e $r_0$ la distanza di equilibrio stabile (punto $P_0$).
Se aumentiamo l'energia del sistema al valore $E_1 > E_{min}$, si ottengono due punti di inversione del moto.
In prima approssimazione $U(r)$ è una parabola per piccole oscillazioni, ma all'aumentare dell'energia, a causa della forma asimmetrica della buca di potenziale, il punto medio dell'oscillazione si sposta verso destra.
Questo significa che l'equilibrio dinamico si raggiunge a una distanza media $r > r_0$; di conseguenza il volume complessivo aumenta poiché le molecole si allontanano mediamente l'una dall'altra. Se l'energia fornita è sufficientemente elevata, le molecole si separano definitivamente (causando ad esempio la fusione della sostanza).

\section{Calore e Quantità di Calore}
\subsection{Stessa Sostanza}

% Due corpi
\begin{figure}[ht]
\centering
\begin{tikzpicture}[scale=1]
\fill[pattern=north east lines, mypen!20] (-0.3,-0.3) rectangle (3.3,1.5);
\draw[thick, mypen] (-0.3,-0.3) rectangle (3.3,1.5);
\draw[thick, fill=cyan!20] (0,0) rectangle (1.4,1.2) node[midway, font=\bfseries] {$m_1, T_1$};
\draw[thick, fill=cyan!40] (1.6,0) rectangle (3,1.2) node[midway, font=\bfseries] {$m_2, T_2$};
\node[above] at (1.5, 1.5) {\footnotesize Contenitore adiabatico (nessun dispersione termica)};
\end{tikzpicture}
\caption{Contatto termico tra due blocchi della stessa sostanza a temperature iniziali differenti all'interno di un involucro adiabatico.}
\end{figure}

Consideriamo due corpi della stessa sostanza con massa e temperatura rispettivamente $m_1, T_1$ e $m_2, T_2$, con $T_1 < T_2$.
Se posti a contatto e isolati dall'esterno, essi raggiungono l'equilibrio termico a una temperatura $T$ intermedia ($T_1 < T < T_2$).

Sperimentalmente si osserva che vale la relazione:
\begin{equation}
m_1 (T - T_1) = m_2 (T_2 - T)
\end{equation}
da cui si ricava che $T$ è la media pesata (con le masse) delle temperature iniziali:
\begin{equation}
T = \frac{m_1 T_1 + m_2 T_2}{m_1 + m_2}
\end{equation}

Generalizzando per $n$ corpi, si ottiene:
\begin{equation}
\sum_{i=1}^n m_i (T - T_i) = 0 \implies T = \frac{\sum m_i T_i}{\sum m_i}
\end{equation}

\vspace{0.3cm}
\subsection{Sostanze Diverse}

% Due corpi
\begin{figure}[ht]
\centering
\begin{tikzpicture}[scale=1]
\fill[pattern=north east lines, mypen!20] (-0.3,-0.3) rectangle (3.3,1.5);
\draw[thick, mypen] (-0.3,-0.3) rectangle (3.3,1.5);
\draw[thick, fill=cyan!20] (0,0) rectangle (1.4,1.2) node[midway, font=\bfseries] {$C_1, T_1$};
\draw[thick, fill=orange!30] (1.6,0) rectangle (3,1.2) node[midway, font=\bfseries] {$C_2, T_2$};
\node[above] at (1.5, 1.5) {\footnotesize Scambio di calore interno: $Q_1 + Q_2 = 0$};
\end{tikzpicture}
\caption{Contatto termico tra due corpi di sostanze diverse (con capacità termiche $C_1$ e $C_2$) posti in isolamento termico rispetto all'ambiente.}
\end{figure}

Consideriamo ora due corpi di sostanze diverse, caratterizzati da massa $m_i$, temperatura $T_i$ e da un coefficiente specifico $c_i$:
\begin{itemize}
    \item $c_i$: \textbf{Calore specifico}, un coefficiente caratteristico della \textit{sostanza} (indipendente dalla massa).
    \item $C_i$: \textbf{Capacità termica}, un coefficiente caratteristico del \textit{corpo} (dipende dalla massa, $C_i = m_i c_i$).
\end{itemize}

Se posti a contatto e isolati dall'esterno, l'equilibrio termico si raggiunge a $T_1 < T < T_2$.
Sperimentalmente si osserva la seguente relazione:
\begin{equation}
m_1 c_1 (T - T_1) = m_2 c_2 (T_2 - T) \iff C_1 (T - T_1) = C_2 (T_2 - T)
\end{equation}

Generalizzando per $n$ corpi, la temperatura finale risulta essere la media pesata con le capacità termiche:
\begin{equation}
\sum_{i=1}^n C_i (T - T_i) = 0 \implies T = \frac{\sum C_i T_i}{\sum C_i}
\end{equation}

\begin{definizione}{Quantità di Calore}
Si definisce \textbf{quantità di calore} $Q_i$ ceduta o assorbita da un corpo:
\begin{equation}
Q_i = C_i (T - T_i) = c_i m_i (T - T_i)
\end{equation}
La convenzione dei segni implica:
\begin{itemize}
    \item Se $T > T_i \implies Q_i > 0$: il corpo \textbf{assorbe} calore.
    \item Se $T < T_i \implies Q_i < 0$: il corpo \textbf{cede} calore.
\end{itemize}
\end{definizione}

\medskip
Considerando l'equazione dell'equilibrio termico $\sum m_i c_i (T - T_i) = 0$ come un sistema omogeneo, abbiamo bisogno di fissare arbitrariamente un calore specifico di riferimento per determinare gli altri. Convenzionalmente si utilizza l'acqua distillata.

\begin{definizione}{Caloria}
La \textbf{chilocaloria} (kcal) è la quantità di calore necessaria per innalzare la temperatura di $1\text{ kg}$ di acqua distillata da $14.5^\circ\text{C}$ a $15.5^\circ\text{C}$ alla pressione di $1\text{ atm}$.
In base a questa definizione, il calore specifico dell'acqua è $c_{\text{acqua}} = 1 \, \frac{\text{kcal}}{\text{kg}^\circ\text{C}}$.
\end{definizione}

\begin{osservazione}{Equivalente meccanico del calore}
La chilocaloria non è una grandezza fisica indipendente, poiché il calore è una forma di energia (la cui unità di misura nel SI è il Joule). Tramite l'esperimento di Joule si è stabilita la conversione:
\begin{equation}
1\text{ kcal} = 4184\text{ J} \implies 1\text{ cal} = 4.184\text{ J}
\end{equation}
(Si nota anche che una frigoria corrisponde a una caloria negativa).
\end{osservazione}

\vspace{0.3cm}
\section{Cambiamenti di Stato}
I cambiamenti di stato (a condizioni standard) sono determinati dalla temperatura:
\begin{itemize}
    \item \textbf{Temperatura di Ebollizione e Liquefazione} ($T_E = T_L$):
    \begin{itemize}
        \item Se $T_{\text{corpo}} > T_E \implies$ gas.
        \item Se $T_{\text{corpo}} < T_E \implies$ liquido.
    \end{itemize}
    \item \textbf{Temperatura di Fusione e Solidificazione} ($T_F = T_S$):
    \begin{itemize}
        \item Se $T_{\text{corpo}} > T_F \implies$ liquido.
        \item Se $T_{\text{corpo}} < T_F \implies$ solido.
    \end{itemize}
    \item \textbf{Temperatura di Brinamento e Sublimazione} ($T_B = T_{\text{sub}}$) [più rari]:
    \begin{itemize}
        \item Se $T_{\text{corpo}} > T_B \implies$ gas.
        \item Se $T_{\text{corpo}} < T_B \implies$ solido.
    \end{itemize}
\end{itemize}

\subsection{Calore Latente}

\begin{definizione}{Calore Latente ($\lambda$)}
Il \textbf{calore latente} è la quantità di calore fornita (o sottratta) a un corpo per passare da uno stato di aggregazione all'altro agendo sui legami intermolecolari, riferita all'unità di massa:
\begin{equation}
    \boxed{ \Delta Q = m \lambda \implies \lambda = \frac{\Delta Q}{m} }
\end{equation}
\end{definizione}

Un tipico esempio è quello del calore latente di fusione $(\Delta Q)_{\text{fusione}}$ o di solidificazione $(\Delta Q)_{\text{solidificazione}}$ dell'acqua ($\text{H}_2\text{O}$).

% Grafico T-Q e T-t
\begin{figure}[ht]
\centering
\begin{tikzpicture}
% Primo grafico T-Q
\draw[->, >=stealth, mypen] (-0.5, 0) -- (4, 0) node[right] {$Q$};
\draw[->, >=stealth, mypen] (0, -1) -- (0, 3) node[above] {$T$};
\draw[mypen] (0, 1) -- (-0.1, 1) node[left] {$0^\circ$};
\draw[thick, mygreen] (-0.5, -0.5) -- (0.5, 1);
\draw[thick, mygreen] (0.5, 1) -- (2.5, 1);
\draw[thick, mygreen] (2.5, 1) -- (3.5, 2.5);
\node at (1.5, 0.7) {$(\Delta Q)_F$};
\draw[dashed, mypen] (1.5, 1) -- (1.5, 0);

% Secondo grafico T-t
\begin{scope}[xshift=6cm]
\draw[->, >=stealth, mypen] (-0.5, 0) -- (4, 0) node[right] {$t$};
\draw[->, >=stealth, mypen] (0, -1) -- (0, 3) node[above] {$T$};
\draw[mypen] (0, 1) -- (-0.1, 1) node[left] {$T_F$};
\draw[thick, blue] (-0.5, -0.5) to[out=45, in=200] (0.5, 1);
\draw[thick, blue] (0.5, 1) -- (2.5, 1);
\draw[thick, blue] (2.5, 1) to[out=20, in=260] (3.5, 2.5);
\node at (1.5, 0.7) {$(\Delta t)_F$};
\draw[dashed, mypen] (1.5, 1) -- (1.5, 0);
\end{scope}
\end{tikzpicture}
\caption{Grafici Temperatura-Calore e Temperatura-tempo per un cambiamento di stato}
\end{figure}

In generale si ha che $(\Delta Q) \propto t$.
Nella fase in cui si fornisce il calore di fusione $(\Delta Q)_{\text{fusione}}$, la temperatura non aumenta e rimane fissa a $T_F = 0^\circ\text{C}$ poiché coesistono ghiaccio e acqua mentre sta avvenendo il passaggio di stato (si tratta di un \textbf{processo isotermo}).
Analizzando i processi inversi (come la solidificazione), il calore scambiato è negativo poiché viene sottratta energia al sistema:
\begin{align}
(\Delta Q)_{\text{fusione}} &= -(\Delta Q)_{\text{solidificazione}} \implies \lambda_F = -\lambda_S = \frac{Q_F}{m} = -\frac{Q_S}{m} \\
(\Delta Q)_{\text{evaporazione}} &= -(\Delta Q)_{\text{liquefazione}} \implies \lambda_E = -\lambda_L = \frac{Q_E}{m} = -\frac{Q_L}{m}
\end{align}

\begin{definizione}{Termostato o Bagno Termico}
Un corpo con massa infinita ($m \to \infty$) e quindi capacità termica infinita ($C \to \infty$) è detto \textbf{termostato}, \textbf{serbatoio}, \textbf{sorgente di calore} o \textbf{bagno termico}. 
La sua caratteristica fondamentale è mantenere la temperatura costante, indipendentemente dal calore scambiato.
\end{definizione}

\medskip
Infatti, se mettiamo in contatto due corpi con capacità termiche $C_1$ e $C_2$, assumendo $C_1 \ll C_2 \to \infty$, la temperatura di equilibrio sarà:
\begin{equation}
T_{\text{equil}} = \frac{C_1 T_1 + C_2 T_2}{C_1 + C_2} \simeq \frac{C_2 T_2}{C_2} = T_2
\end{equation}

\vspace{0.3cm}
\section{Propagazione del Calore}
Il calore si propaga in tre modi differenti:
\begin{itemize}
    \item \textbf{Convezione} (Legge di Newton)
    \item \textbf{Conduzione} (Legge di Fourier)
    \item \textbf{Irraggiamento} (Legge di Stefan-Boltzmann)
\end{itemize}

\subsection{Convezione}
\begin{definizione}{Convezione}
La \textbf{convezione} è un meccanismo di trasmissione del calore caratterizzato da un macroscopico \textbf{trasporto di materia}, tipico dei fluidi (liquidi e gas), in cui il movimento di masse fluide trasferisce energia termica da una regione all'altra.
\end{definizione}

\begin{figure}[ht]
\centering
\begin{tikzpicture}[scale=1]
\fill[cyan!10] (0,0) rectangle (2.5, 2.2);
\draw[thick, mypen] (0,0) rectangle (2.5, 2.2);
\draw[->, >=stealth, thick, notered] (0.6, 0.4) to[out=90, in=180] (1.25, 1.8) to[out=0, in=90] (1.9, 0.4);
\draw[->, >=stealth, thick, blue] (1.9, 0.4) to[out=270, in=0] (1.25, 0.1) to[out=180, in=270] (0.6, 0.4);
\draw[->, >=stealth, thick, notered, line width=1.5pt] (1.25, -0.6) -- (1.25, 0) node[midway, right] {$Q$};
\fill[pattern=north east lines, notered!30] (-0.2, -0.6) rectangle (2.7, 0);
\draw[thick, notered] (-0.2, 0) -- (2.7, 0);
\node[font=\bfseries, mypen] at (3.5, 1.1) {\parbox{3cm}{\centering Moti convettivi nel fluido}};
\end{tikzpicture}
\caption{Moti convettivi in un fluido riscaldato dal basso: il fluido caldo, meno denso, sale verso l'alto ed è sostituito dal fluido freddo che scende, instaurando un ciclo termico continuo con trasporto di materia.}
\end{figure}

All'interno di un fluido riscaldato dal basso si generano delle correnti che ciclicamente si scaldano (salendo verso l'alto per via della minore densità) e si raffreddano (scendendo nuovamente).

\begin{definizione}{Legge di Newton (empirica)}
La quantità di calore scambiata nell'unità di tempo fra l'unità di superficie di un corpo e il fluido circostante è proporzionale alla differenza di temperatura:
\begin{equation}
\boxed{Q = h S t \Delta T}
\end{equation}
dove $h$ è il \textbf{coefficiente di conduttività esterna} (determinato empiricamente).
\end{definizione}

\subsection{Conduzione}
\begin{definizione}{Conduzione}
La \textbf{conduzione} è un meccanismo di propagazione del calore in un mezzo materiale che avviene tramite urti e interazioni a livello molecolare \textbf{senza alcun trasporto di materia} (a differenza della convezione).
\end{definizione}

\begin{figure}[ht]
\centering
\begin{tikzpicture}[scale=1]
\fill[pattern=north east lines, notered!25] (0,0) rectangle (0.6, 2);
\draw[thick, mypen, fill=orange!10] (0.6, 0) rectangle (3.4, 2);
\fill[pattern=north west lines, cyan!25] (3.4, 0) rectangle (4, 2);
\draw[->, >=stealth, thick, mypen] (0.6, -0.5) -- (3.8, -0.5) node[right] {$x$};
\node[font=\bfseries] at (0.3, 2.3) {Sorgente Calda ($T_2$)};
\node[font=\bfseries] at (3.7, 2.3) {Sorgente Fredda ($T_1$)};
\draw[<->, >=stealth, thick, mypen] (1.5, 1.5) -- (2.5, 1.5) node[midway, above] {$\Delta x$};
\draw[->, >=stealth, thick, notered, line width=1.2pt] (1.2, 0.8) -- (2.8, 0.8) node[midway, below, font=\bfseries] {Flusso $Q$};
\node at (0.6, -0.8) {0};
\node at (3.4, -0.8) {$L$};
\end{tikzpicture}
\caption{Conduzione termica lungo un conduttore cilindrico isolato lateralmente di lunghezza $L$, posto tra due serbatoi a temperature fisse $T_2 > T_1$.}
\end{figure}

La conduzione termica si verifica fra due solidi a contatto oppure all'interno di uno stesso solido.
Consideriamo due zone a temperature $T_2$ e $T_1$ con $T_2 > T_1$. Fissiamo un asse $x$ con origine sulla superficie che divide il corpo dalla zona a $T_2$.

\begin{figure}[ht]
\centering
\begin{tikzpicture}[scale=1]
\draw[->, >=stealth, thick, mypen] (0,0) -- (3.5,0) node[right] {$x$};
\draw[->, >=stealth, thick, mypen] (0,0) -- (0,2.8) node[above] {$T(x)$};
\draw[dashed, mypen] (0,2.2) -- (2.8,2.2);
\draw[dashed, mypen] (0,0.6) -- (2.8,0.6);
\draw[dashed, mypen] (2.8,0) -- (2.8,0.6);
\node[left, font=\bfseries, notered] at (0, 2.2) {$T_2$};
\node[left, font=\bfseries, blue] at (0, 0.6) {$T_1$};
\node[below] at (2.8, 0) {$L$};
\draw[thick, mygreen, line width=1.5pt] (0,2.2) -- (2.8, 0.6);
\end{tikzpicture}
\caption{Andamento lineare della temperatura $T(x)$ all'interno di un conduttore omogeneo in regime stazionario, caratterizzato da un gradiente termico costante $\frac{dT}{dx} = \frac{T_1-T_2}{L} < 0$.}
\end{figure}

Spostandoci dalla zona a $T_2$ verso quella a $T_1$, la temperatura del corpo diminuisce gradualmente fino a raggiungere l'estremità opposta; di conseguenza esiste un \textbf{gradiente di temperatura} negativo: $\frac{dT}{dx} < 0$.

\begin{definizione}{Legge di Fourier}
La quantità di calore scambiata nell'unità di tempo fra due solidi (o due strati contigui) è proporzionale alla superficie di contatto e al gradiente termico. La formula sperimentale per un flusso su una lunghezza $L$ in regime stazionario è:
\begin{equation}
\boxed{Q = \frac{k S t \Delta T}{L}}
\end{equation}
dove $k$ è il \textbf{coefficiente di conduttività interna} (dipendente dal materiale).
\begin{itemize}
    \item Se $k$ è \textbf{elevato}, il materiale è un buon conduttore termico.
    \item Se $k$ è \textbf{ridotto}, il materiale è un buon isolante termico.
\end{itemize}
\end{definizione}

\medskip
Il ragionamento alla base della formula infinitesima di Fourier considera le seguenti proporzionalità per il flusso di calore $\delta Q$:
\begin{itemize}
    \item $\delta Q \propto S$ (superficie di contatto)
    \item $\delta Q \propto -\frac{dT}{dx}$ (gradiente termico, il segno meno indica il flusso verso temperature inferiori)
    \item $\delta Q \propto k$ (conduttività)
\end{itemize}

\begin{figure}[ht]
\centering
\begin{tikzpicture}[scale=1.2]
\draw[thick, fill=orange!20] (0,0) rectangle (0.6, 1.5);
\draw[dashed, mypen] (0,0) -- (0,-0.4);
\draw[dashed, mypen] (0.6,0) -- (0.6,-0.4);
\draw[<->, >=stealth, thick, mypen] (0,-0.3) -- (0.6,-0.3) node[midway, below] {$dx$};
\node[left, font=\bfseries] at (0, 0.75) {$T + dT$};
\node[right, font=\bfseries] at (0.6, 0.75) {$T$};
\draw[->, >=stealth, thick, notered, line width=1.5pt] (-0.5, 0.75) -- (1.2, 0.75) node[right, font=\bfseries] {$\delta Q$};
\end{tikzpicture}
\caption{Bilancio termico su uno strato infinitesimo di spessore $dx$: il calore fluisce nel verso delle temperature decrescenti ($\delta Q \propto -dT/dx$).}
\end{figure}

Da queste considerazioni si ricava la forma differenziale della Legge di Fourier:
\begin{equation}
\boxed{\frac{\delta Q}{dt} = -k S \frac{dT}{dx}}
\end{equation}
da cui, integrando e supponendo $k$ uniforme, si ottiene la formula precedente.

Se invece abbiamo una conduzione fra materiali differenti in serie (ovvero $k$ non uniforme):

\begin{figure}[ht]
\centering
\begin{tikzpicture}[scale=1]
\fill[pattern=north east lines, notered!20] (-0.4,0) rectangle (0,2.5);
\draw[thick, fill=orange!10] (0,0) rectangle (1.2,2.5) node[midway] {$k_1$};
\draw[thick, fill=yellow!15] (1.2,0) rectangle (2.6,2.5) node[midway] {$k_2$};
\draw[thick, fill=green!10] (2.6,0) rectangle (4,2.5) node[midway] {$k_3$};
\fill[pattern=north west lines, cyan!20] (4,0) rectangle (4.4,2.5);

\draw[<->, >=stealth] (0,-0.3) -- (1.2,-0.3) node[midway, below] {$d_1$};
\draw[<->, >=stealth] (1.2,-0.3) -- (2.6,-0.3) node[midway, below] {$d_2$};
\draw[<->, >=stealth] (2.6,-0.3) -- (4,-0.3) node[midway, below] {$d_3$};

\node[left, font=\bfseries, notered] at (0, 2.8) {$T_2$};
\node[font=\bfseries] at (1.2, 2.8) {$T_A$};
\node[font=\bfseries] at (2.6, 2.8) {$T_B$};
\node[right, font=\bfseries, blue] at (4, 2.8) {$T_1$};

\draw[->, >=stealth, thick, notered, line width=1.5pt] (0.8, -1) -- (3.2, -1) node[midway, below, font=\bfseries] {Flusso stazionario $Q_{\text{TOT}}$};
\end{tikzpicture}
\caption{Conduzione termica attraverso una parete composta da tre strati di materiale differente in serie ($T_2 > T_A > T_B > T_1$). In regime stazionario il flusso di calore $Q_{\text{TOT}}$ è uniforme attraverso ogni strato.}
\end{figure}

\begin{align}
Q_1 &= \frac{k_1 S (T_2 - T_A) t}{d_1} \\
Q_2 &= \frac{k_2 S (T_A - T_B) t}{d_2} \\
Q_3 &= \frac{k_3 S (T_B - T_1) t}{d_3}
\end{align}

In condizioni stazionarie il calore totale fluito attraverso ciascuno strato deve essere lo stesso: $Q_1 = Q_2 = Q_3 = Q_{\text{TOT}}$.
Esprimendo le differenze di temperatura e sommandole:
\begin{equation}
(T_2 - T_A) + (T_A - T_B) + (T_B - T_1) = T_2 - T_1
\end{equation}
da cui si ottiene la relazione per il calore totale ceduto da $T_2$ a $T_1$:
\begin{equation}
\boxed{Q_{\text{TOT}} = \frac{S (T_2 - T_1) t}{\frac{d_1}{k_1} + \frac{d_2}{k_2} + \frac{d_3}{k_3}}}
\end{equation}
Questa relazione è alla base del principio dei doppi vetri: la quantità rilevante per l'isolamento è il rapporto $\frac{d_i}{k_i}$.

\vspace{0.3cm}
\subsection{Irraggiamento}
\begin{definizione}{Irraggiamento}
L'\textbf{irraggiamento} è la trasmissione di energia termica mediante emissione e propagazione di \textbf{onde elettromagnetiche} da parte di qualsiasi corpo a temperatura maggiore dello zero assoluto ($T > 0\text{ K}$). A differenza della conduzione e della convezione, avviene anche nel vuoto assoluto senza necessità di alcun mezzo materiale in interposizione.
\end{definizione}

\begin{definizione}{Legge di Stefan-Boltzmann}
La quantità di calore irradiata da un corpo nero nell'unità di tempo è proporzionale alla quarta potenza della sua temperatura assoluta:
\begin{equation}
\boxed{Q = S t \sigma T^4}
\end{equation}
dove $\sigma = 5.670 \cdot 10^{-8} \, \frac{\text{W}}{\text{m}^2 \text{K}^4}$ è la \textbf{costante di Stefan-Boltzmann}.
(Per i corpi non neri la costante dipende dall'emissività della sostanza).
\end{definizione}

\begin{figure}[ht]
\centering
\begin{tikzpicture}[scale=1.1]
% Cavità del corpo nero (sezione di sfera con foro)
\fill[gray!20] (0,0) circle (1.5);
\fill[pattern=north east lines, mypen!40] (0,0) circle (1.5);
\fill[white] (0,0) circle (1.2);
\fill[white] (1.1,-0.3) rectangle (1.6,0.3);
\draw[thick, mypen] (1.3,0.3) arc[start angle=12, end angle=348, radius=1.35];

% Raggio incidente e riflessioni interne
\draw[->, >=stealth, thick, notered, line width=1.2pt] (2.5, 0.1) -- (1.3, 0.1) node[midway, above] {Radiazione incidente};
\draw[->, >=stealth, thick, notered] (1.3, 0.1) -- (-0.8, 0.7);
\draw[->, >=stealth, thick, notered] (-0.8, 0.7) -- (-0.5, -0.9);
\draw[->, >=stealth, thick, notered] (-0.5, -0.9) -- (0.7, -0.6);
\draw[->, >=stealth, thick, notered] (0.7, -0.6) -- (0.2, 0.9);
\draw[->, >=stealth, thick, notered] (0.2, 0.9) -- (-1.1, -0.2);

\node[font=\bfseries, mypen] at (0, -1.8) {Cavità isoterma (Corpo Nero ideale)};
\end{tikzpicture}
\caption{Modello fisico di Corpo Nero realizzato tramite una cavità con una piccola apertura: ogni raggio di luce che entra subisce innumerevoli riflessioni interne venendo completamente assorbito dalle pareti, senza possibilità di uscirne.}
\end{figure}

\vspace{0.3cm}
\section{Concetti Fondamentali della Termodinamica}

\begin{definizione}{Sistema Termodinamico e Ambiente}
Un \textbf{sistema termodinamico} è una porzione di materia (con composizione chimica nota) separata dall'ambiente circostante tramite una superficie di controllo. 
Se le proprietà chimico-fisiche sono le stesse in tutto il sistema, esso è detto \textbf{omogeneo}. 
Si distinguono tre tipologie di sistemi:
\begin{itemize}
    \item \textbf{Aperti}: scambiano sia materia che energia con l'ambiente.
    \item \textbf{Chiusi}: scambiano solo energia, non materia.
    \item \textbf{Isolati}: non scambiano né materia né energia.
\end{itemize}
\end{definizione}

\begin{figure}[ht]
\centering
\begin{tikzpicture}
    % Sistema generico
    \draw[thick, fill=cyan!10, mypen] (0,0) circle (2);
    \node at (0,0) {Sistema};
    % Frecce di scambio
    \draw[->, >=stealth, thick, notered] (2.5, 1) -- (1.5, 0.5) node[midway, above] {Calore $Q$};
    \draw[<-, >=stealth, thick, mypen] (2.5, -1) -- (1.5, -0.5) node[midway, below] {Lavoro $L$};
\end{tikzpicture}
\caption{Rappresentazione di un sistema termodinamico aperto/chiuso.}
\end{figure}

Lo \textbf{stato del sistema} è definito da un insieme di grandezze macroscopiche dette \textbf{variabili di stato} (pressione $p$, volume $V$, temperatura $T$). Esse sono legate tra loro dall'equazione di stato $f(p, V, T) = 0$. Ad esempio, per un gas perfetto vale $p V = n R T$.

\begin{definizione}{Equilibrio Termodinamico}
Un sistema si trova in stato di \textbf{equilibrio termodinamico} se soddisfa contemporaneamente tre condizioni:
\begin{enumerate}
    \item \textbf{Equilibrio Meccanico}: La risultante delle forze interne e delle forze esterne è nulla ($\sum \vec{F}_{\text{int}} = 0$, $\sum \vec{F}_{\text{ext}} = 0$).
    \item \textbf{Equilibrio Termico}: La temperatura $T$ è uniforme in tutto il sistema.
    \item \textbf{Equilibrio Chimico}: Assenza di reazioni chimiche macroscopiche e di spostamento netto di materia.
\end{enumerate}
In questo stato la situazione è stazionaria, è descrivibile da poche variabili di stato ed è indipendente dalla storia pregressa del sistema.
\end{definizione}

\vspace{0.3cm}
\subsection{Trasformazioni Termodinamiche}
Una trasformazione termodinamica è un processo durante il quale alcune variabili di stato variano da una situazione iniziale $A$ (di equilibrio) a una situazione finale $B$ (di equilibrio).
Se lo stato finale coincide con lo stato iniziale ($A \equiv B$), la trasformazione è detta \textbf{ciclica}.

Inoltre, se la trasformazione soddisfa le seguenti condizioni:
\begin{itemize}
    \item È una successione continua di stati di equilibrio infinitamente vicini (così che le variabili di stato siano definite in ogni istante).
    \item È infinitamente lenta (il sistema ha il tempo per riequilibrarsi).
    \item Le cause che la generano sono infinitamente piccole.
\end{itemize}
allora la trasformazione è definita \textbf{reversibile} (cioè invertibile per il sistema e per l'ambiente senza lasciare alcuna traccia). Altrimenti, la trasformazione è \textbf{irreversibile} (ogni processo reale caratterizzato da attrito meccanico, turbolenza o gradienti finiti di temperatura è irreversibile).

\begin{osservazione}{Criterio Globale di Reversibilità e Ciclicità}
Una distinzione cruciale tra processi reversibili e irreversibili risiede negli effetti sull'ambiente:
\begin{itemize}
    \item \textbf{Processo Reversibile}: Effettuando la trasformazione diretta $A \to B$ e successivamente la trasformazione inversa $B \to A$, sia il \textbf{sistema} che l'\textbf{ambiente} tornano esattamente nei rispettivi stati iniziali, senza che nell'universo rimanga alcuna traccia residua del ciclo compiuto.
    \item \textbf{Processo Irreversibile}: Anche se è sempre possibile riportare il \textbf{sistema} allo stato iniziale $A$ (facendogli compiere un ciclo chiuso per cui $\Delta U_{\text{sist}} = 0$), l'\textbf{ambiente} subirà inevitabilmente una modificazione permanente (ad esempio un dispendio netto di lavoro trasformato in calore dissipato nel termostato freddo).
\end{itemize}
\end{osservazione}

\begin{osservazione}{Accoppiamento Termico e Catena Infinita di Termostati}
Consideriamo un corpo di capacità termica $C$ inizialmente a temperatura $T_1$. Se lo poniamo bruscamente a contatto con un singolo termostato a temperatura $T_2 > T_1$, il corpo assorbe una quantità di calore $Q = C(T_2 - T_1)$. Se successivamente lo raffreddiamo ponendolo a contatto con un termostato a $T_1$, esso cede il calore $Q' = -C(T_2 - T_1)$. 
Per il corpo il processo è ciclico ($\Delta U = 0$). Tuttavia, per l'ambiente si è verificato un trasferimento diretto di calore dalla sorgente calda $T_2$ alla sorgente fredda $T_1$ senza alcuna produzione di lavoro, un processo fortemente \textbf{irreversibile}.
L'unico modo concettuale per rendere lo scambio termico \textbf{quasi-reversibile} consiste nell'accoppiare il corpo con una \textbf{catena infinita di termostati ausiliari} a temperature $T_1, T_1+\delta T, T_1+2\delta T, \dots, T_2$, con differenze infinitesime $\delta T \to 0$, in modo che in ogni istante il corpo e il termostato in contatto siano in equilibrio termico approssimato.
\end{osservazione}

\begin{osservazione}{Rappresentazione nel Piano di Clapeyron ($p, V$)}
Nel diagramma di Clapeyron, una trasformazione può essere rappresentata da una \textbf{curva continua} e ben definita \textbf{solo se è reversibile}, poiché solo in tal caso ogni punto intermedio corrisponde a uno stato di equilibrio macroscopico con pressione $p$ e temperatura $T$ uniformi in tutto il volume. 
Per una trasformazione \textbf{irreversibile}, durante il processo intermedio il sistema non è in equilibrio e le variabili di stato non hanno valori uniformi o definiti; pertanto, \textbf{la curva $p(V)$ non esiste}. Graficamente, una trasformazione irreversibile tra due stati di equilibrio $A$ e $B$ si indica convenzionalmente con una \textbf{linea spezzata o tratteggiata}.
\end{osservazione}

\begin{figure}[ht]
\centering
\begin{tikzpicture}[scale=1]
% Espansione
\fill[pattern=north east lines, mypen!20] (0,0) rectangle (0.6, 1.6);
\draw[thick, mypen] (0,0) rectangle (0.6, 1.6);
\draw[thick, mypen, fill=cyan!10] (0.6, 0) rectangle (2.2, 1.6);
\draw[thick, fill=gray!30] (2.2, 0) rectangle (2.4, 1.6);
\draw[->, >=stealth, thick, mypen, line width=1.2pt] (2.4, 0.8) -- (3.6, 0.8) node[right, font=\bfseries] {$F$};
\draw[<-, >=stealth, thick, notered] (0.6, 1.3) -- (1.8, 1.3) node[right] {$pS$};
\draw[<-, >=stealth, thick, mygreen] (0.6, 0.3) -- (1.8, 0.3) node[right] {$F_A$};
\node[above, font=\bfseries] at (2.8, 1.2) {Vuoto};
\node[below, font=\bfseries, blue] at (1.4, -0.2) {Espansione ($dV > 0$)};
\node[right] at (4, 0.8) {
    \parbox{7cm}{\small
    \textbf{Equilibrio dinamico:}\\
    $pS = F + F_A \implies F = pS - F_A < pS$ \\
    \textit{La forza motrice esterna richiesta è minore rispetto al caso ideale senza attrito.}
    }
};

% Compressione
\begin{scope}[yshift=-3.8cm]
\fill[pattern=north east lines, mypen!20] (0,0) rectangle (0.6, 1.6);
\draw[thick, mypen] (0,0) rectangle (0.6, 1.6);
\draw[thick, mypen, fill=orange!10] (0.6, 0) rectangle (2.2, 1.6);
\draw[thick, fill=gray!30] (2.2, 0) rectangle (2.4, 1.6);
\draw[<-, >=stealth, thick, mypen, line width=1.2pt] (2.4, 0.8) -- (3.6, 0.8) node[right, font=\bfseries] {$F$};
\draw[->, >=stealth, thick, notered] (0.6, 1.3) -- (1.8, 1.3) node[right] {$pS$};
\draw[<-, >=stealth, thick, mygreen] (0.6, 0.3) -- (1.8, 0.3) node[right] {$F_A$};
\node[above, font=\bfseries] at (2.8, 1.2) {Vuoto};
\node[below, font=\bfseries, notered] at (1.4, -0.2) {Compressione ($dV < 0$)};
\node[right] at (4, 0.8) {
    \parbox{7cm}{\small
    \textbf{Equilibrio dinamico:}\\
    $pS + F_A = F \implies F = pS + F_A > pS$ \\
    \textit{La forza motrice esterna richiesta è maggiore rispetto al caso ideale senza attrito.}
    }
};
\end{scope}
\end{tikzpicture}
\caption{Confronto meccanico tra espansione e compressione di un gas in un cilindro con pistone in presenza di forza d'attrito viscoso $F_A$: l'istèresi meccanica dimostra l'irreversibilità dei processi reali.}
\end{figure}

Di conseguenza, $F_{\text{ESP}} \neq F_{\text{COM}}$, quindi in presenza di attrito il percorso di andata è sempre diverso dal percorso di ritorno, confermando l'irreversibilità del processo reale.
\section{Primo Principio della Termodinamica}

\begin{definizione}{Calore, Lavoro ed Energia Interna}
Il Primo Principio della Termodinamica è un'estensione del principio di conservazione dell'energia. L'energia interna $U$ di un sistema è una \textbf{funzione di stato}; la sua variazione $\Delta U$ tra due stati $A$ e $B$ dipende solo dagli stati estremi e non dal percorso seguito:
\begin{equation}
    \Delta U = U_B - U_A = Q - L
\end{equation}
dove:
\begin{itemize}
    \item $Q$ è il calore scambiato dal sistema (positivo se assorbito, negativo se ceduto).
    \item $L$ è il lavoro compiuto (positivo se compiuto dal sistema sull'ambiente, negativo se subito).
\end{itemize}
\end{definizione}

In forma differenziale si scrive:
\begin{equation}
    dU = \delta Q - \delta L
\end{equation}

\begin{osservazione}{Differenziali Inesatti ($\delta Q$ e $\delta L$)}
Si nota l'uso del simbolo $\delta$ (o đ) per calore e lavoro, in contrapposizione al differenziale esatto $dU$. Questo indica che $Q$ e $L$ \textbf{non sono funzioni di stato}: le quantità infinitesime $\delta Q$ e $\delta L$ non sono i differenziali di una funzione del solo punto (stato), ma dipendono dal cammino seguito lungo la trasformazione. L'integrale $\int_A^B \delta L = L_{AB}$ rappresenta il lavoro lungo una specifica traiettoria, mentre $\int_A^B dU = U_B - U_A$ è lo stesso per qualunque cammino.
\end{osservazione}

\subsection{Lavoro delle Forze di Pressione (Trasformazioni Reversibili e Irreversibili)}
Consideriamo un fluido racchiuso all'interno di un cilindro di sezione di area $S$ munito di un pistone mobile. Se il pistone si sposta di un tratto infinitesimo $dx$, la forza esterna che si oppone all'espansione è $F_e = p_e S$, dove $p_e$ è la pressione esterna. Il lavoro infinitesimo compiuto dal sistema contro l'ambiente esterno è:
\begin{equation}
    \delta L = F_e \, dx = p_e S \, dx = p_e \, dV
\end{equation}
(Il lavoro compiuto dalle forze esterne sul sistema è di segno opposto, $\delta L_e = -p_e dV$).

\begin{itemize}
    \item \textbf{Trasformazione Reversibile}: In ogni istante il sistema è in equilibrio meccanico con l'ambiente, per cui la pressione interna del gas è uniforme e bilancia esattamente quella esterna ($p = p_e$). Il lavoro totale per una variazione finita di volume da $V_A$ a $V_B$ è dato dall'integrale:
    \begin{equation}
        \boxed{ L_{\text{rev}} = \int_{V_A}^{V_B} p(V) \, dV }
    \end{equation}
    Nel piano di Clapeyron ($p, V$), questo integrale rappresenta geometricamente l'\textbf{area sottesa dalla curva} della trasformazione.
    \item \textbf{Trasformazione Irreversibile}: Durante una trasformazione irreversibile non esiste una pressione $p$ interna uniforme definita, per cui l'integrale $\int p \, dV$ perde di significato rispetto alle coordinate del sistema. Il lavoro può essere calcolato \textbf{solo valutando le forze esterne}, ovvero $L_{\text{irr}} = \int p_e \, dV$. Se ad esempio un gas si espande liberamente nel vuoto ($p_e = 0$), il lavoro compiuto è rigorosamente nullo ($L = 0$), nonostante il volume aumenti da $V_A$ a $V_B$.
\end{itemize}

\begin{figure}[ht]
\centering
\begin{tikzpicture}
    % Asse p-V
    \draw[->, >=stealth, thick, mypen] (0,0) -- (0,4) node[above] {$p$};
    \draw[->, >=stealth, thick, mypen] (0,0) -- (5,0) node[right] {$V$};
    
    % Area sotto la curva per il lavoro
    \fill[blue!10, opacity=0.5] (1,0) -- (1,3) to[out=-60, in=150] (4,1) -- (4,0) -- cycle;
    
    % Trasformazione
    \draw[thick, blue, ->, >=stealth] (1,3) to[out=-60, in=150] (4,1);
    
    % Punti A e B
    \filldraw[mypen] (1,3) circle (2pt) node[above right] {$A$};
    \filldraw[mypen] (4,1) circle (2pt) node[above right] {$B$};
    
    \node at (2.5, 1) {$L = \int p \, dV$};
\end{tikzpicture}
\caption{Lavoro termodinamico di una trasformazione reversibile come area sottesa nel piano di Clapeyron ($p-V$).}
\end{figure}

\vspace{0.3cm}
\section{Calorimetria e Capacità Termica}

\begin{definizione}{Calore Specifico Molare}
Oltre al calore specifico per unità di massa ($c = \frac{1}{m}\frac{\delta Q}{dT}$), per i gas è più utile riferirsi al numero di moli $n$. Si definiscono due \textbf{calori specifici molari}:
\begin{itemize}
    \item $c_v$: calore specifico molare a volume costante.
    \item $c_p$: calore specifico molare a pressione costante.
\end{itemize}
Per solidi e liquidi la dilatazione termica è trascurabile ($dV \simeq 0$), sicché il lavoro di espansione è nullo e le due capacità termiche coincidono ($c_p \simeq c_v \simeq c$). Per i gas, invece, la differenza è notevole ed è regolata dalla \textbf{relazione di Mayer}:
\begin{equation}
    c_p - c_v = R
\end{equation}
\end{definizione}

\begin{definizione}{Calori Specifici dei Gas Perfetti, Gradi di Libertà ed Energia Interna}
In base al teorema di equipartizione dell'energia della meccanica statistica, ogni grado di libertà molecolare contribuisce per $\frac{1}{2}R$ al calore specifico molare a volume costante $c_v$. Definendo il rapporto dei calori specifici $\gamma = \frac{c_p}{c_v} > 1$, si ha:
\begin{itemize}
    \item \textbf{Gas Monoatomici} (3 gradi di libertà traslazionali):
    \begin{equation}
        c_v = \frac{3}{2}R, \quad c_p = \frac{5}{2}R, \quad \gamma = \frac{5}{3} \simeq 1.67
    \end{equation}
    \item \textbf{Gas Biatomici} (3 traslazionali + 2 rotazionali a temperatura ambiente):
    \begin{equation}
        c_v = \frac{5}{2}R, \quad c_p = \frac{7}{2}R, \quad \gamma = \frac{7}{5} = 1.40
    \end{equation}
    \item \textbf{Gas Poliatomici} (3 traslazionali + 3 rotazionali):
    \begin{equation}
        c_v = 3R, \quad c_p = 4R, \quad \gamma = \frac{4}{3} \simeq 1.33
    \end{equation}
\end{itemize}
Per un gas perfetto, l'energia interna $U$ è funzione esclusiva della temperatura (Legge di Joule). La variazione di energia interna tra due stati a temperature $T_A$ e $T_B$ vale pertanto:
\begin{equation}
    \boxed{ \Delta U = n c_v \Delta T = n c_v (T_B - T_A) }
\end{equation}
Questa relazione fondamentale è valida per \textbf{qualunque trasformazione} (isocora, isobara, isoterma, adiabatica o persino irreversibile), poiché $U$ è una funzione di stato.
\end{definizione}

\begin{definizione}{Legge di Dulong e Petit per i Solidi}
Per i solidi cristallini a temperatura ambiente (al di sopra della temperatura caratteristica di Debye), ogni atomo oscilla attorno alla propria posizione di equilibrio reticolare comportandosi come un oscillatore armonico tridimensionale (3 gradi di libertà cinetici e 3 potenziali). Il calore specifico molare approssima un valore universale indipendente dalla natura chimica della sostanza:
\begin{equation}
    \boxed{ C_{\text{mol}} \simeq 6 \cdot \frac{1}{2}R = 3R \simeq 24.9 \, \frac{\text{J}}{\text{mol} \cdot \text{K}} \simeq 6 \, \frac{\text{cal}}{\text{mol} \cdot \text{K}} }
\end{equation}
\end{definizione}

\begin{figure}[ht]
\centering
\begin{tikzpicture}
    % Calorimetro
    \draw[thick, mypen] (0,0) -- (0,3) -- (3,3) -- (3,0) -- cycle;
    \draw[thick, mypen] (0.2,0.2) -- (0.2,2.8) -- (2.8,2.8) -- (2.8,0.2) -- cycle;
    
    % Termometro
    \draw[thick, mypen, fill=notered] (1.5, 0.5) circle (0.15);
    \draw[thick, mypen] (1.4, 0.5) -- (1.4, 4) -- (1.6, 4) -- (1.6, 0.5);
    \node at (1.5, 4.2) {$T$};
    
    % Acqua
    \fill[cyan!20] (0.2, 0.2) rectangle (2.8, 2);
    
    \node at (1.5, -0.5) {Calorimetro};
\end{tikzpicture}
\caption{Schema semplificato di un calorimetro.}
\end{figure}

\vspace{0.3cm}
\section{Tipi di Trasformazioni Termodinamiche}
Un sistema può passare da uno stato iniziale a uno finale attraverso diverse trasformazioni notevoli:
\begin{itemize}
    \item \textbf{Isoterma} ($T = \text{costante}$): 
    \begin{equation}
        p V = \text{cost.} \implies L = n R T \ln\left(\frac{V_B}{V_A}\right), \quad \Delta U = 0
    \end{equation}
    
    \item \textbf{Isocora} ($V = \text{costante}$):
    \begin{equation}
        \frac{p}{T} = \text{cost.} \implies L = 0, \quad Q = \Delta U = n c_v \Delta T
    \end{equation}
    
    \item \textbf{Isobara} ($p = \text{costante}$):
    \begin{equation}
        \frac{V}{T} = \text{cost.} \implies L = p \Delta V, \quad Q = n c_p \Delta T
    \end{equation}
    
    \item \textbf{Adiabatica} ($Q = 0$):
    \begin{equation}
        p V^\gamma = \text{cost.} \quad \text{con} \quad \gamma = \frac{c_p}{c_v}
    \end{equation}
    Il lavoro si esprime come $L = -\Delta U = -n c_v \Delta T$.
\end{itemize}

\begin{definizione}{Espansione Libera di Joule (Adiabatica Irreversibile)}
L'\textbf{espansione libera di Joule} è un processo in cui un gas perfetto, contenuto in un recipiente rigidamente isolato termicamente ($Q = 0$), si espande nel vuoto all'apertura di una valvola di comunicazione. Poiché l'espansione avviene contro una pressione esterna nulla ($p_e = 0$), il lavoro compiuto è nullo ($L = 0$). Dal Primo Principio della Termodinamica segue:
\begin{equation}
    \Delta U = Q - L = 0 \implies U_B = U_A
\end{equation}
Essendo l'energia interna di un gas perfetto funzione esclusiva della temperatura, si conclude che $T_B = T_A$ ($\Delta T = 0$). 
\textbf{Attenzione}: Sebbene la temperatura iniziale e finale coincidano, questa \textbf{non è una trasformazione isoterma reversibile}. Durante il moto turbolento nel vuoto, il gas attraversa stati di non-equilibrio in cui pressione e temperatura non sono definite uniformemente. Il processo è \textbf{fortemente irreversibile} e comporta un aumento di entropia dell'universo $\Delta S = n R \ln(V_B / V_A) > 0$. Nel diagramma di Clapeyron i punti estremi $A$ e $B$ giacciono sulla medesima iperbole isoterma, ma non possono essere congiunti da alcuna curva continua.
\end{definizione}

\begin{figure}[ht]
\centering
\begin{tikzpicture}[scale=1.1]
    % Griglia di fondo leggera
    \draw[step=1cm, gray!15, very thin] (0,0) grid (5.2,5.2);

    % Assi
    \draw[->, >=stealth, thick, mypen] (0,0) -- (0,5.2) node[above] {$p$};
    \draw[->, >=stealth, thick, mypen] (0,0) -- (5.2,0) node[right] {$V$};
    
    % Punto di intersezione comune A(2, 3) -> p*V = 6 per isoterma; p*V^1.4 = 7.917 per adiabatica
    % Isocora (V = 2)
    \draw[thick, mygreen] (2, 0.8) -- (2, 4.8) node[above] {\small Isocora};
    
    % Isobara (p = 3)
    \draw[thick, orange] (0.8, 3) -- (4.8, 3) node[right] {\small Isobara};
    
    % Isoterma (p = 6/V)
    \draw[thick, blue, domain=1.2:4.8, samples=50] plot (\x, {6/\x}) node[right] {\small Isoterma};
    
    % Adiabatica (p = 7.917 / V^1.4)
    \draw[thick, notered, domain=1.35:4.8, samples=50] plot (\x, {7.917/(\x^1.4)}) node[right] {\small Adiabatica};
    
    % Punto di intersezione e proiezioni
    \draw[dashed, gray] (2,0) node[below] {$V_0$} -- (2,3) -- (0,3) node[left] {$p_0$};
    \filldraw[mypen] (2,3) circle (2pt) node[above right=1pt] {$A$};
\end{tikzpicture}
\caption{Confronto tra le quattro trasformazioni fondamentali nel piano $p-V$, intersecantesi nel medesimo stato iniziale $A(V_0, p_0)$. Si noti come, durante un'espansione ($V > V_0$), la curva adiabatica scenda più rapidamente rispetto all'isoterma a causa dell'assenza di apporto di calore ($\gamma > 1$).}
\end{figure}
\vspace{0.3cm}
\subsection{Derivazione dell'Adiabatica Reversibile}
Nella trasformazione adiabatica non c'è scambio di calore ($Q = 0$), da cui per il primo principio si ottiene $L = -\Delta U$.
In forma infinitesima per un gas perfetto reversibile si ha:
\begin{equation}
    \delta Q = \delta L + dU = p \, dV + n c_v \, dT = 0
\end{equation}

Dall'equazione di stato $pV = nRT$ differenziando si ricava:
\begin{equation}
    dp \, V + p \, dV = n R \, dT \implies n \, dT = \frac{V}{R} dp + \frac{p}{R} dV
\end{equation}

Sostituendo il termine $n \, dT$ nell'equazione del primo principio:
\begin{equation}
    p \, dV + c_v \left( \frac{V}{R} dp + \frac{p}{R} dV \right) = 0
\end{equation}

Raccogliendo i termini, considerando che $R = c_p - c_v$ e dividendo tutto per $c_v/R$:
\begin{equation}
    p \left( \frac{c_v + R}{c_v} \right) dV + V \, dp = 0 \implies p \frac{c_p}{c_v} dV + V \, dp = 0
\end{equation}

Posto il coefficiente adiabatico $\gamma = \frac{c_p}{c_v}$, si ricava:
\begin{equation}
    \gamma \frac{dV}{V} = - \frac{dp}{p}
\end{equation}

Essendo la trasformazione reversibile, è possibile integrare lungo la curva passante per gli stati $A$ e $B$:
\begin{equation}
    \gamma \int_{V_A}^{V_B} \frac{dV}{V} = -\int_{p_A}^{p_B} \frac{dp}{p} \implies \gamma \ln\left(\frac{V_B}{V_A}\right) + \ln\left(\frac{p_B}{p_A}\right) = 0
\end{equation}
Applicando le proprietà dei logaritmi si giunge infine alle \textbf{equazioni di Poisson} per l'adiabatica reversibile:
\begin{equation}
    \boxed{p V^\gamma = \text{cost.}} \quad \quad \boxed{T V^{\gamma-1} = \text{cost.}} \quad \quad \boxed{p^{1-\gamma} T^\gamma = \text{cost.}}
\end{equation}

\vspace{0.3cm}
\section{Secondo Principio e Macchine Termiche}

\begin{definizione}{Enunciati del Secondo Principio}
Esistono due formulazioni storiche equivalenti del Secondo Principio della Termodinamica:
\begin{itemize}
    \item \textbf{Enunciato di Kelvin-Planck}: È impossibile realizzare una trasformazione il cui unico risultato sia quello di assorbire calore da una sola sorgente e trasformarlo integralmente in lavoro.
    \item \textbf{Enunciato di Clausius}: È impossibile realizzare una trasformazione il cui unico risultato sia il passaggio di calore da un corpo a temperatura minore a uno a temperatura maggiore.
\end{itemize}
\end{definizione}

\medskip
Una \textbf{macchina termica} opera ciclicamente tra due sorgenti a temperature $T_c$ (calda) e $T_f$ (fredda). Il suo \textbf{rendimento} $\eta$ è definito come il rapporto tra il lavoro prodotto $L$ e il calore assorbito dalla sorgente calda $Q_c$:
\begin{equation}
    \eta = \frac{L}{Q_c} = \frac{Q_c - |Q_f|}{Q_c} = 1 - \frac{|Q_f|}{Q_c}
\end{equation}

Per il ciclo reversibile di \textbf{Carnot}, che rappresenta il rendimento massimo teorico operando tra due temperature, vale:
\begin{equation}
    \eta_{\text{Carnot}} = 1 - \frac{T_f}{T_c}
\end{equation}

\begin{figure}[ht]
\centering
\begin{tikzpicture}
    % Sorgente Calda
    \draw[thick, mypen, fill=notered!30] (2,4) rectangle (4,5);
    \node at (3,4.5) {Sorgente $T_c$};
    
    % Macchina Termica
    \draw[thick, mypen, fill=gray!20] (3,2.5) circle (0.8);
    \node at (3,2.5) {M.T.};
    
    % Sorgente Fredda
    \draw[thick, mypen, fill=cyan!30] (2,0) rectangle (4,1);
    \node at (3,0.5) {Sorgente $T_f$};
    
    % Frecce
    \draw[->, >=stealth, thick, notered] (3,4) -- (3,3.3) node[midway, right] {$Q_c$};
    \draw[->, >=stealth, thick, cyan] (3,1.7) -- (3,1) node[midway, right] {$Q_f$};
    \draw[->, >=stealth, thick, mypen] (3.8, 2.5) -- (5, 2.5) node[midway, above] {$L$};
\end{tikzpicture}
\caption{Schema di principio di una macchina termica.}
\end{figure}

\vspace{0.3cm}
\subsection{Espansione Libera (Trasformazione Irreversibile)}

Se il gas espande nel vuoto non compie lavoro macroscopico contro forze esterne e, qualora il recipiente sia adiabatico, non scambia calore.

\begin{figure}[ht]
\centering
\begin{tikzpicture}[scale=1.1]
    % Pareti adiabatiche esterne
    \fill[pattern=north east lines, mypen!25] (-0.3,-0.3) rectangle (5.5,1.8);
    \fill[white] (0,0) rectangle (5.2,1.5);
    \draw[thick, mypen] (0,0) rectangle (5.2,1.5);
    
    % Prima scatola A (Stato A)
    \draw[thick, fill=cyan!20] (0,0) rectangle (2.3,1.5);
    \fill[pattern=north west lines, cyan!40] (0,0) rectangle (2.3,1.5);
    \draw[thick, mypen, line width=1.5pt] (2.3,0) -- (2.3,1.5);
    \node[font=\bfseries] at (1.15, 0.75) {Gas ($V_A, T_A$)};
    \node[font=\bfseries, gray] at (3.75, 0.75) {Vuoto ($p=0$)};
    \node[below, font=\bfseries] at (2.6, -0.4) {Stato A (Iniziale): Setto chiuso};

    % Seconda scatola B (Stato B)
    \begin{scope}[yshift=-3cm]
    \fill[pattern=north east lines, mypen!25] (-0.3,-0.3) rectangle (5.5,1.8);
    \fill[white] (0,0) rectangle (5.2,1.5);
    \draw[thick, mypen] (0,0) rectangle (5.2,1.5);
    \draw[thick, fill=cyan!15] (0,0) rectangle (5.2,1.5);
    \draw[dashed, mypen] (2.3,0) -- (2.3,1.5);
    \node[font=\bfseries] at (2.6, 0.75) {Gas espanso all'equilibrio ($V_B > V_A, T_B = T_A$)};
    \node[below, font=\bfseries] at (2.6, -0.4) {Stato B (Finale): Setto rimosso};
    \end{scope}
    
    \draw[->, >=stealth, thick, mypen, line width=1.5pt] (2.6, -0.9) -- (2.6, -1.3) node[midway, right] {Espansione libera};
\end{tikzpicture}
\caption{Esperimento di espansione libera di Joule nel vuoto all'interno di un contenitore rigido e adiabatico: poiché $L=0$ e $Q=0$, per il Primo Principio si ha $\Delta U = 0$ e, per un gas perfetto, la temperatura rimane invariata ($T_B = T_A$).}
\end{figure}

Poiché il gas si espande nel vuoto, la pressione esterna è nulla ($p_e = 0$) e non viene compiuto lavoro ($L = 0$). Inoltre, essendo il recipiente isolato termicamente, non vi è scambio di calore ($Q = 0$). Per il Primo Principio della Termodinamica si ha dunque:
\begin{equation}
    \Delta U = Q - L = 0
\end{equation}
Per un gas perfetto, l'energia interna dipende unicamente dalla temperatura ($U = U(T)$), quindi $\Delta U = 0 \implies \Delta T = 0$. La trasformazione tra lo stato iniziale $A$ e finale $B$ giunge alla stessa temperatura ($T_A = T_B$), comportandosi complessivamente come una trasformazione \textbf{isoterma} che segue la legge di Boyle:
\begin{equation}
    p_A V_A = p_B V_B
\end{equation}

\vspace{0.3cm}
\subsection{Confronto delle Pendenze tra Adiabatica e Isoterma}
Nel piano di Clapeyron ($p-V$), una curva adiabatica reversibile risulta sempre più ripida di una curva isoterma passante per il medesimo punto.

\begin{figure}[ht]
\centering
\begin{minipage}{0.42\textwidth}
\centering
\begin{tikzpicture}[scale=1]
    % Griglia
    \draw[step=1cm, gray!15, very thin] (0,0) grid (4.2,4.2);
    % Assi
    \draw[->, >=stealth, thick, mypen] (0,0) -- (0,4.2) node[left] {$p$};
    \draw[->, >=stealth, thick, mypen] (0,0) -- (4.2,0) node[below] {$V$};
    
    % Curve passanti per P(1.8, 2.2): p*V = 3.96; p*V^1.4 = 5.019
    \draw[thick, blue, domain=1.0:3.8, samples=40] plot (\x, {3.96/\x}) node[right] {\small Isoterma};
    \draw[thick, notered, domain=1.15:3.6, samples=40] plot (\x, {5.019/(\x^1.4)}) node[right] {\small Adiabatica};
    
    % Tangente Isoterma (pendenza -1.22)
    \draw[thin, blue!70, dashed, domain=1.1:2.7] plot (\x, {4.4 - 1.22*\x});
    
    % Tangente Adiabatica (pendenza -1.71)
    \draw[thin, notered!70, dashed, domain=1.25:2.5] plot (\x, {5.28 - 1.71*\x});
    
    % Punto di contatto
    \filldraw[mypen] (1.8, 2.2) circle (1.8pt) node[above right=1pt] {$P$};
    \draw[dashed, gray] (1.8,0) node[below] {$V_0$} -- (1.8,2.2) -- (0,2.2) node[left] {$p_0$};
\end{tikzpicture}
\end{minipage}\hfill
\begin{minipage}{0.54\textwidth}
Differenziando l'equazione dell'adiabatica reversibile $pV^\gamma = \text{cost.}$:
\begin{align*}
    d(pV^\gamma) = 0 &\implies V^\gamma dp + \gamma V^{\gamma-1} p\,dV = 0 \\
    &\implies \left( \frac{dp}{dV} \right)_{\text{adiab}} = - \gamma \frac{p}{V}
\end{align*}
Confrontando con la pendenza dell'isoterma ($pV = \text{cost.} \implies \gamma = 1$):
\begin{align}
    \left( \frac{dp}{dV} \right)_{\text{isot}} &= - \frac{p}{V} \\
    \left( \frac{dp}{dV} \right)_{\text{adiab}} &= - \gamma \frac{p}{V} = \gamma \left( \frac{dp}{dV} \right)_{\text{isot}}
\end{align}
Poiché $\gamma = \frac{c_p}{c_v} > 1$, la pendenza (in modulo) della curva adiabatica nel punto di intersezione $P(V_0, p_0)$ è maggiore di quella dell'isoterma esattamente di un fattore $\gamma$.
\end{minipage}
\caption{Confronto analitico e geometrico delle pendenze tra una curva isoterma (blu) e un'adiabatica reversibile (rossa) passanti per il medesimo stato $P(V_0, p_0)$. Le rette tratteggiate indicano le rispettive tangenti in $P$.}
\end{figure}

\vspace{0.3cm}
\subsection{Macchine Frigorifere}
A differenza della macchina termica, una \textbf{macchina frigorifera} (o pompa di calore) opera assorbendo un lavoro $L < 0$ dall'ambiente per estrarre calore $Q_f > 0$ da una sorgente fredda a temperatura $T_f$ e cederlo come calore $Q_c < 0$ a una sorgente calda a temperatura $T_c$.

\begin{figure}[ht]
\centering
\begin{tikzpicture}
    % Sorgente Calda
    \draw[thick, mypen, fill=notered!30] (2,4) rectangle (4,5);
    \node at (3,4.5) {Sorgente $T_c$};
    
    % Macchina Frigorifera
    \draw[thick, mypen, fill=gray!20] (3,2.5) circle (0.8);
    \node at (3,2.5) {M.F.};
    
    % Sorgente Fredda
    \draw[thick, mypen, fill=cyan!30] (2,0) rectangle (4,1);
    \node at (3,0.5) {Sorgente $T_f$};
    
    % Frecce
    \draw[->, >=stealth, thick, notered] (3,3.3) -- (3,4) node[midway, right] {$|Q_c|$};
    \draw[->, >=stealth, thick, cyan] (3,1) -- (3,1.7) node[midway, right] {$Q_f$};
    \draw[<-, >=stealth, thick, mypen] (3.8, 2.5) -- (5, 2.5) node[midway, above] {$|L|$};
\end{tikzpicture}
\caption{Schema di principio di una macchina frigorifera.}
\end{figure}

Essendo un ciclo, la variazione di energia interna è nulla, per cui il lavoro scambiato equivale alla somma algebrica dei calori: $|L| = |Q_c| - Q_f$.
L'efficienza di questa macchina è descritta dal \textbf{Coefficiente di Prestazione (COP)}, definito come il rapporto tra l'effetto utile (il calore estratto dalla sorgente fredda) e il costo in termini di lavoro assorbito $|L|$:
\begin{equation}
    \text{COP} = \frac{Q_f}{|L|} = \frac{Q_f}{|Q_c| - Q_f}
\end{equation}

\vspace{0.3cm}
\section{Ciclo di Carnot e Teorema di Carnot}

\begin{figure}[ht]
\centering
\begin{minipage}{0.48\textwidth}
\centering
\begin{tikzpicture}[scale=1.1]
    % Griglia
    \draw[step=1cm, gray!15, very thin] (0,0) grid (4.5,4.5);
    % Assi
    \draw[->, >=stealth, thick, mypen] (0,0) -- (0,4.6) node[left] {$p$};
    \draw[->, >=stealth, thick, mypen] (0,0) -- (4.6,0) node[below] {$V$};
    
    % Curve analitiche del ciclo
    % A->B: Isoterma T2 (p = 4.8/V)
    \draw[thick, blue, domain=1.2:3.0, samples=30] plot (\x, {4.8/\x});
    \node[right, blue] at (2.9, 1.8) {$T_2$};
    
    % B->C: Adiabatica espansione (p = 7.45/V^1.4)
    \draw[thick, notered, domain=3.0:4.0, samples=30] plot (\x, {7.45/(\x^1.4)});
    
    % C->D: Isoterma T1 (p = 2.4/V)
    \draw[thick, blue, domain=1.6:4.0, samples=30] plot (\x, {2.4/\x});
    \node[right, blue] at (3.9, 0.4) {$T_1$};
    
    % D->A: Adiabatica compressione (p = 5.13/V^1.4)
    \draw[thick, notered, domain=1.2:1.6, samples=30] plot (\x, {5.13/(\x^1.4)});
    
    % Riempimento area ciclo (Lavoro L)
    \fill[cyan!15, opacity=0.7, domain=1.2:3.0, variable=\x] (1.2,4.0) -- plot (\x, {4.8/\x}) -- (3.0,1.6) -- (3.0,0.8) -- plot[domain=3.0:1.6] (\x, {2.4/\x}) -- (1.6,1.5) -- plot[domain=1.6:1.2] (\x, {5.13/(\x^1.4)}) -- cycle;
    
    % Vertici
    \filldraw[mypen] (1.2, 4.0) circle (1.8pt) node[above right] {$A$};
    \filldraw[mypen] (3.0, 1.6) circle (1.8pt) node[above right] {$B$};
    \filldraw[mypen] (4.0, 0.6) circle (1.8pt) node[below right] {$C$};
    \filldraw[mypen] (1.6, 1.5) circle (1.8pt) node[below left] {$D$};
    
    \node[mypen, font=\bfseries] at (2.4, 2.0) {$L > 0$};
\end{tikzpicture}
\end{minipage}\hfill
\begin{minipage}{0.50\textwidth}
Il \textbf{Ciclo di Carnot} è un ciclo termodinamico reversibile operante tra due sorgenti a temperatura fissa $T_2 > T_1$. Si compone di quattro fasi alternate:
\begin{itemize}
    \item \textbf{$A \to B$}: Espansione isoterma a temperatura $T_2$ \\
    ($Q_2 = Q_{AB} = L_{AB} > 0$)
    \item \textbf{$B \to C$}: Espansione adiabatica reversibile \\
    ($Q = 0 \implies L_{BC} = -\Delta U > 0$)
    \item \textbf{$C \to D$}: Compressione isoterma a temperatura $T_1$ \\
    ($Q_1 = Q_{CD} = L_{CD} < 0$)
    \item \textbf{$D \to A$}: Compressione adiabatica reversibile \\
    ($Q = 0 \implies L_{DA} = -\Delta U < 0$)
\end{itemize}
\end{minipage}
\caption{Rappresentazione nel piano di Clapeyron del ciclo reversibile di Carnot. L'area racchiusa dal ciclo (azzurra) rappresenta il lavoro utile $L = Q_2 - |Q_1|$ prodotto dalla macchina in ogni ciclo.}
\end{figure}

\vspace{0.3cm}
Essendo il ciclo chiuso, $\Delta U_{\text{tot}} = 0$. Il calore assorbito è $Q_2 = Q_{AB}$ e quello ceduto $Q_1 = Q_{CD}$. Dalle leggi delle adiabatiche si dimostra che $\frac{V_B}{V_A} = \frac{V_C}{V_D}$, per cui il rendimento si semplifica in:
\begin{equation}
    \eta_C = \frac{L}{Q_2} = 1 - \frac{|Q_1|}{Q_2} = 1 - \frac{T_1}{T_2}
\end{equation}

\begin{definizione}{Teorema di Carnot}
Date due sorgenti di calore a temperature $T_1$ e $T_2$ (con $T_2 > T_1$):
\begin{enumerate}
    \item Tutte le macchine termiche \textbf{reversibili} operanti tra le due sorgenti hanno lo stesso rendimento $\eta_C = 1 - T_1/T_2$.
    \item Nessuna macchina termica \textbf{irreversibile} può avere un rendimento maggiore: $\eta_i \le \eta_C$.
\end{enumerate}
\end{definizione}

\vspace{0.3cm}
\begin{figure}[ht]
\centering
\begin{minipage}{0.42\textwidth}
\centering
\begin{tikzpicture}[scale=1.05]
    % Sorgente CALDA T2
    \draw[thick, mypen, fill=notered!20, rounded corners=2pt] (0, 3) rectangle (3, 3.8) node[midway, font=\bfseries] {Sorgente Calda $T_2$};
    
    % Macchina Mc (Carnot in senso inverso - Frigorifero)
    \draw[thick, mypen, fill=white] (0.75, 1.5) circle (0.5) node[font=\bfseries] {$M_C$};
    % Calore ceduto a T2 (freccia verso l'alto da Mc a T2)
    \draw[->, >=stealth, thick, notered] (0.75, 2) -- (0.75, 3) node[left, midway] {$Q_{2C}$};
    % Calore assorbito da T1 (freccia verso l'alto da T1 a Mc)
    \draw[->, >=stealth, thick, cyan] (0.75, 0) -- (0.75, 1) node[left, midway] {$Q_{1C}$};
    % Lavoro assorbito da Mc (freccia entrante in Mc da sinistra)
    \draw[->, >=stealth, thick, mypen] (-0.3, 1.5) node[left] {$L_C$} -- (0.25, 1.5);
    
    % Macchina M (Generica - Motore termico)
    \draw[thick, mypen, fill=white] (2.25, 1.5) circle (0.5) node[font=\bfseries] {$M$};
    % Calore assorbito da T2 (freccia verso il basso da T2 a M)
    \draw[->, >=stealth, thick, notered] (2.25, 3) -- (2.25, 2) node[right, midway] {$Q_2$};
    % Calore ceduto a T1 (freccia verso il basso da M a T1)
    \draw[->, >=stealth, thick, cyan] (2.25, 1) -- (2.25, 0) node[right, midway] {$Q_1$};
    % Lavoro prodotto da M (freccia uscente da M verso destra)
    \draw[->, >=stealth, thick, mypen] (2.75, 1.5) -- (3.3, 1.5) node[right] {$L$};
    
    % Sorgente FREDDA T1
    \draw[thick, mypen, fill=cyan!20, rounded corners=2pt] (0, -0.8) rectangle (3, 0) node[midway, font=\bfseries] {Sorgente Fredda $T_1$};
\end{tikzpicture}
\end{minipage}\hfill
\begin{minipage}{0.54\textwidth}
\textbf{Dimostrazione del Teorema di Carnot}:\\
Si accoppia una macchina termica generica $M$ a una macchina reversibile di Carnot $M_C$ operante in senso inverso (come frigorifero). La macchina $M$ produce un lavoro $L > 0$ estraendo calore dalla sorgente $T_2$, mentre $M_C$ assorbe lavoro $L_C < 0$ per rimettere calore nella medesima sorgente.

Regolando le dimensioni dei cicli in modo che il saldo termico con la sorgente calda sia nullo ($Q_2 + Q_{2C} = 0$), il sistema combinato $M + M_C$ scambia calore unicamente con la sorgente fredda $T_1$.
\end{minipage}
\caption{Schema di accoppiamento per la dimostrazione del Teorema di Carnot via assurdo secondo l'enunciato di Kelvin-Planck.}
\end{figure}

\vspace{0.3cm}
Per il Secondo Principio (enunciato di Kelvin-Planck), il lavoro totale prodotto non può essere positivo (non si può estrarre calore da un'unica sorgente $T_1$ e produrre lavoro netto). Quindi $L_{tot} \le 0 \implies L - |L_C| \le 0$. Da questo limite termodinamico deriva la disuguaglianza sui rendimenti:
\begin{itemize}
    \item Se $M$ è \textbf{irreversibile}: si dimostra che $\eta \le \eta_C$.
    \item Se $M$ è \textbf{reversibile}: il ciclo $M + M_C$ è reversibile e può essere invertito, portando a dedurre che $\eta = \eta_C$.
\end{itemize}

\vspace{0.3cm}
\noindent $\implies$ La macchina reversibile è la più efficiente ($Q \leftrightarrow L$) perché, a parità di calore assorbito, cede la minima quantità di calore alla sorgente fredda.

\begin{osservazione}{Irreversibilità e Moto Perpetuo}
Il primo principio non fa distinzioni tra processi reversibili e irreversibili, tuttavia l'esperienza quotidiana mostra che i processi spontanei (come il calore che fluisce da un corpo caldo a uno freddo, o l'energia dissipata per attrito) sono sempre irreversibili. Questo stabilisce una preferenza temporale (la \textit{freccia del tempo}).

Il \textbf{Secondo Principio} codifica queste osservazioni, vietando:
\begin{itemize}
    \item \textbf{Macchine anti-Clausius}: frigoriferi che estraggono calore da una sorgente fredda verso una calda senza richiedere alcun lavoro esterno in ingresso.
    \item \textbf{Moto perpetuo di seconda specie (anti-Kelvin-Planck)}: macchine termiche in grado di trasformare ciclicamente calore in lavoro scambiando con un'unica sorgente termica.
\end{itemize}
Di conseguenza, l'unica possibilità per un dispositivo che compie un ciclo interagendo con \textbf{una sola sorgente termica (ciclo monotermo)} è quella di subire lavoro per trasformarlo integralmente in calore ($L \le 0$ e $Q \le 0$).
\end{osservazione}

\vspace{0.5cm}
\subsection{Esempio: Ciclo monotermo per un gas perfetto}
Consideriamo un ciclo termodinamico compiuto da $n$ moli di un gas perfetto che scambia calore con una sola sorgente a temperatura $T_1$. Il ciclo è costituito da tre trasformazioni:
\begin{itemize}
    \item \textbf{A $\to$ B (Espansione isoterma reversibile a $T_1$)}: Il gas si espande da $V_A$ a $V_B > V_A$ scambiando calore con l'unico termostato a temperatura $T_1$. Il lavoro compiuto è:
    \begin{equation}
        L_{AB} = n R T_1 \ln\left(\frac{V_B}{V_A}\right)
    \end{equation}
    \item \textbf{B $\to$ C (Compressione adiabatica reversibile)}: Il gas viene compresso fino al volume iniziale $V_C = V_A$, portando la temperatura a $T_2 \neq T_1$. Senza scambio di calore ($Q_{BC}=0$), il lavoro compiuto è pari all'opposto della variazione di energia interna:
    \begin{equation}
        L_{BC} = -\Delta U_{BC} = -n c_v (T_2 - T_1)
    \end{equation}
    \item \textbf{C $\to$ A (Raffreddamento isocoro irreversibile)}: Mantenendo il volume costante a $V_A$, il gas viene rimescolato e riportato alla temperatura iniziale $T_1$ cedendo calore. Essendo isocora, non si compie lavoro: $L_{CA} = 0$. (Nota: l'accoppiamento termico tra una temperatura $T_2$ e l'unico termostato a $T_1$ rende la fase irreversibile).
\end{itemize}

\begin{figure}[ht]
\centering
\begin{minipage}{0.48\textwidth}
\centering
\begin{tikzpicture}[scale=1.1]
    % Griglia
    \draw[step=1cm, gray!15, very thin] (0,0) grid (3.8,3.2);
    % Assi
    \draw[->, >=stealth, thick, mypen] (0,0) -- (0,3.3) node[left] {$p$};
    \draw[->, >=stealth, thick, mypen] (0,0) -- (3.8,0) node[below] {$V$};
    
    % Vertici analitici (VA = VC = 1.0, VB = 3.0)
    % Isoterma T1: p = 1.5/V -> A(1.0, 1.5), B(3.0, 0.5)
    % Adiabatica T1->T2: p = 2.33 / V^1.4 -> B(3.0, 0.5), C(1.0, 2.33)
    
    % A->B Isoterma
    \draw[thick, blue, domain=1.0:3.0, samples=30] plot (\x, {1.5/\x});
    % B->C Adiabatica
    \draw[thick, notered, domain=1.0:3.0, samples=30] plot (\x, {2.33/(\x^1.4)});
    % C->A Isocora
    \draw[thick, mygreen, ->, >=stealth] (1.0, 2.33) -- (1.0, 1.5);
    
    % Frecce di direzione sulle curve
    \draw[->, >=stealth, thick, blue] (1.8, {1.5/1.8}) -- (1.9, {1.5/1.9});
    \draw[<-, >=stealth, thick, notered] (1.7, {2.33/(1.7^1.4)}) -- (1.8, {2.33/(1.8^1.4)});
    
    % Punti
    \filldraw[mypen] (1.0, 1.5) circle (1.6pt) node[left] {$A$};
    \filldraw[mypen] (3.0, 0.5) circle (1.6pt) node[right] {$B$};
    \filldraw[mypen] (1.0, 2.33) circle (1.6pt) node[left] {$C$};
    
    \node[below, blue] at (2.0, 0.7) {\small Isoterma ($T_1$)};
    \node[above right, notered] at (1.6, 1.4) {\small Adiabatica};
    \node[left, mygreen] at (1.0, 1.9) {\small Isocora};
\end{tikzpicture}
\end{minipage}\hfill
\begin{minipage}{0.48\textwidth}
Il lavoro totale del ciclo è la somma dei lavori delle singole trasformazioni:
\begin{equation}
    L_{\text{tot}} = L_{AB} + L_{BC} + L_{CA} = n R T_1 \ln\left(\frac{V_B}{V_A}\right) - n c_v (T_2 - T_1)
\end{equation}
Essendo $L_{CA} = 0$, il segno del lavoro totale dipende dal confronto tra il lavoro positivo di espansione isoterma e quello negativo assorbito nella compressione adiabatica.
\end{minipage}
\caption{Ciclo monotermo costituito da un'isoterma reversibile ($A \to B$), un'adiabatica reversibile ($B \to C$) e un'isocora irreversibile ($C \to A$).}
\end{figure}

Per l'adiabatica reversibile $B \to C$ vale la relazione tra temperature e volumi:
\begin{equation}
    \frac{T_2}{T_1} = \left(\frac{V_B}{V_C}\right)^{\gamma-1} = \left(\frac{V_B}{V_A}\right)^{\gamma-1}
\end{equation}
Passando ai logaritmi e ricordando che $\gamma - 1 = \frac{c_p - c_v}{c_v} = \frac{R}{c_v}$, si ottiene:
\begin{equation}
    \ln\left(\frac{T_2}{T_1}\right) = (\gamma - 1) \ln\left(\frac{V_B}{V_A}\right) = \frac{R}{c_v} \ln\left(\frac{V_B}{V_A}\right) \implies n R \ln\left(\frac{V_B}{V_A}\right) = n c_v \ln\left(\frac{T_2}{T_1}\right)
\end{equation}
Sostituendo questa espressione nel lavoro totale $L_{\text{tot}}$:
\begin{align}
    L_{\text{tot}} &= n c_v T_1 \ln\left(\frac{T_2}{T_1}\right) - n c_v (T_2 - T_1) \nonumber \\
    &= n c_v T_1 \left[ \ln\left(\frac{T_2}{T_1}\right) - \left(\frac{T_2}{T_1} - 1\right) \right]
\end{align}
Per la disuguaglianza fondamentale dei logaritmi, per qualsiasi $x > 0$ con $x \neq 1$ vale sempre $\ln(x) < x - 1$, ossia $\ln(x) - (x-1) < 0$. Posto $x = \frac{T_2}{T_1}$, risulta rigorosamente:
\begin{equation}
    \boxed{L_{\text{tot}} < 0}
\end{equation}
Questo dimostra matematicamente che in un ciclo monotermo irreversibile il lavoro totale prodotto è necessariamente negativo: il sistema subisce lavoro netto dall'ambiente per dissiparlo in calore, in perfetto accordo con il Secondo Principio della Termodinamica (se infatti fosse reversibile o producesse lavoro positivo $L > 0$, violerebbe l'enunciato di Kelvin-Planck).




\subsection{Equivalenza degli Enunciati di Kelvin-Planck e di Clausius}
Per dimostrare l'equivalenza logica dei due enunciati ($K \iff C$), è sufficiente far vedere che se uno dei due è falso, allora necessariamente lo è anche l'altro. Procediamo con due dimostrazioni per assurdo:

\vspace{0.3cm}
\noindent \textbf{1. Dimostrazione $K \implies C$ (se vale Kelvin-Planck, deve valere Clausius)}: \\
Supponiamo per assurdo che l'enunciato di Kelvin-Planck sia vero e quello di Clausius sia falso.
Esistono allora una \textbf{macchina anti-Clausius} $C$ (in grado di trasferire calore $Q_1 > 0$ da una sorgente fredda $T_1$ a una calda $T_2$ senza apporto di lavoro esterno) e una normale macchina termica di Carnot $M_C$ operante tra le stesse due sorgenti.

\begin{figure}[ht]
\centering
\begin{minipage}{0.40\textwidth}
\centering
\begin{tikzpicture}[scale=0.95]
    % T2
    \draw[thick, mypen, fill=notered!20, rounded corners=2pt] (0, 3) rectangle (3, 3.8) node[midway, font=\bfseries] {Sorgente $T_2 > T_1$};
    % T1
    \draw[thick, mypen, fill=cyan!20, rounded corners=2pt] (0, -0.8) rectangle (3, 0) node[midway, font=\bfseries] {Sorgente $T_1$};
    
    % Macchina anti-Clausius C
    \draw[thick, mypen, fill=white] (0.75, 1.5) circle (0.5) node[font=\bfseries] {$C$};
    \draw[->, >=stealth, thick, cyan] (0.75, 0) -- (0.75, 1) node[midway, left] {$Q_1$};
    \draw[->, >=stealth, thick, notered] (0.75, 2) -- (0.75, 3) node[midway, left] {$Q_1$};
    
    % Macchina termica di Carnot Mc (Motore normale)
    \draw[thick, mypen, fill=white] (2.25, 1.5) circle (0.5) node[font=\bfseries] {$M_C$};
    \draw[->, >=stealth, thick, notered] (2.25, 3) -- (2.25, 2) node[midway, right] {$Q_2$};
    \draw[->, >=stealth, thick, cyan] (2.25, 1) -- (2.25, 0) node[midway, right] {$Q_1$};
    \draw[->, >=stealth, thick, mypen] (2.75, 1.5) -- (3.7, 1.5) node[right, font=\bfseries] {$L$};
\end{tikzpicture}
\end{minipage}\hfill
\begin{minipage}{0.56\textwidth}
Accoppiando le due macchine in modo che scambino la stessa quantità di calore $Q_1$ con la sorgente fredda $T_1$, analizziamo i bilanci termici per il sistema combinato:
\begin{itemize}
    \item La sorgente fredda $T_1$ cede calore $Q_1$ alla macchina $C$ e riceve la medesima quantità $Q_1$ dalla macchina $M_C \implies Q_{\text{netto}, T_1} = 0$.
    \item La sorgente calda $T_2$ riceve calore $Q_1$ dalla macchina $C$ e cede calore $Q_2 > Q_1$ alla macchina $M_C \implies Q_{\text{netto}, T_2} = Q_2 - Q_1 > 0$.
    \item Il lavoro totale prodotto dal sistema combinato, essendo $\Delta U_{\text{tot}} = 0$ per un ciclo, vale:
    \begin{equation}
        L = Q_2 - Q_1 > 0
    \end{equation}
\end{itemize}
\end{minipage}
\caption{Violazione dell'enunciato di Kelvin-Planck indotta dall'accoppiamento con una ipotetica macchina anti-Clausius $C$.}
\end{figure}

Il sistema combinato si comporta quindi come una macchina termica che compie un ciclo interagendo con \textbf{un'unica sorgente di calore} ($T_2$, dato che la sorgente $T_1$ non subisce variazioni nette di calore) e produce un lavoro interamente positivo ($L > 0$). Ciò costituisce una diretta violazione dell'enunciato di Kelvin-Planck, contraddicendo l'ipotesi iniziale. Dunque:
\begin{equation}
    K \text{ vero} \implies C \text{ vero}
\end{equation}

\vspace{0.3cm}
\noindent \textbf{2. Dimostrazione $C \implies K$ (se vale Clausius, deve valere Kelvin-Planck)}: \\
Supponiamo per assurdo che l'enunciato di Clausius sia vero e quello di Kelvin-Planck sia falso.
Esiste allora una \textbf{macchina anti-Kelvin-Planck} $K$ capace di prelevare calore $Q > 0$ da una sola sorgente $T_1$ e trasformarlo integralmente in lavoro meccanico $L = Q$.

\begin{figure}[ht]
\centering
\begin{minipage}{0.46\textwidth}
\centering
\begin{tikzpicture}[scale=0.9]
    % T1
    \draw[thick, mypen, fill=cyan!20, rounded corners=2pt] (0, 3) rectangle (1.8, 4) node[midway, font=\bfseries] {Sorgente $T_1$};
    
    % Macchina anti-Kelvin K
    \draw[thick, mypen, fill=white] (0.9, 1.5) circle (0.5) node[font=\bfseries] {$K$};
    \draw[->, >=stealth, thick, cyan] (0.9, 3) -- (0.9, 2) node[midway, right] {$Q$};
    \draw[->, >=stealth, thick, mypen] (1.4, 1.5) -- (2.6, 1.5) node[midway, above] {$L$};
    
    % Attrito / Dissipatore
    \draw[thick, dashed, mypen, fill=orange!15, rounded corners=2pt] (2.6, 0.8) rectangle (4.2, 2.2) node[midway, align=center, font=\small] {Dissipatore\\(attrito)};
    \draw[->, >=stealth, thick, notered] (4.2, 1.5) -- (5.2, 1.5) node[midway, above] {$Q$};
    
    % T2
    \draw[thick, mypen, fill=notered!20, rounded corners=2pt] (5.2, 0.8) rectangle (7, 2.2) node[midway, align=center, font=\bfseries] {Sorgente\\$T_2 > T_1$};
\end{tikzpicture}
\end{minipage}\hfill
\begin{minipage}{0.50\textwidth}
Utilizziamo ora il lavoro $L$ prodotto dalla macchina $K$ per generare calore tramite attrito (o resistenza elettrica) all'interno di un corpo o di un bagno termico a temperatura superiore $T_2 > T_1$.
\begin{itemize}
    \item La sorgente fredda $T_1$ cede calore netto $Q$.
    \item Il lavoro meccanico $L = Q$ viene dissipato integralmente in calore all'interno della sorgente calda $T_2$.
    \item Nessun altro lavoro netto entra o esce dall'ambiente esterno.
\end{itemize}
\end{minipage}
\caption{Violazione dell'enunciato di Clausius indotta dall'accoppiamento con una ipotetica macchina anti-Kelvin-Planck $K$.}
\end{figure}

Il risultato complessivo del sistema accoppiato è il puro trasferimento di una quantità di calore $Q > 0$ da un corpo freddo ($T_1$) a un corpo caldo ($T_2$) senza alcun apporto netto di lavoro dall'esterno. Ciò viola l'enunciato di Clausius, contraddicendo l'ipotesi iniziale. Dunque:
\begin{equation}
    C \text{ vero} \implies K \text{ vero}
\end{equation}
Rimane così provata la totale equivalenza fisica tra le due formulazioni del Secondo Principio ($K \iff C$).

\section{Temperatura Termodinamica Assoluta}
Dal Teorema di Carnot sappiamo che il rendimento di una macchina reversibile operante tra due sorgenti dipende unicamente dalle loro temperature. Da ciò si deduce che il rapporto tra i calori scambiati in un ciclo di Carnot dipende solo dalle temperature empiriche $t_1, t_2$ delle due sorgenti:
\begin{equation*}
    \frac{|Q_1|}{|Q_2|} = f(t_1, t_2)
\end{equation*}

Accoppiando opportunamente due macchine di Carnot operanti rispettivamente tra $(t_2, t_0)$ e $(t_0, t_1)$, Lord Kelvin dimostrò che tale funzione universale deve essere necessariamente fattorizzabile nella forma $f(t_1, t_2) = \frac{\theta(t_1)}{\theta(t_2)}$.

Si può quindi definire la \textbf{temperatura termodinamica assoluta} come $T \equiv \theta(t)$. Questa nuova scala è intrinseca e del tutto \textbf{indipendente dal fluido termodinamico} (al contrario dei termometri reali). Fortunatamente, si dimostra che essa coincide con la temperatura definita tramite il termometro a gas perfetto:
\begin{equation}
    \boxed{\frac{|Q_1|}{|Q_2|} = \frac{T_1}{T_2}} \implies \eta_C = 1 - \frac{T_1}{T_2}
\end{equation}

\vspace{0.3cm}
\begin{osservazione}{Conseguenze sulla macchina termica}
Poiché le temperature termodinamiche $T$ sono per definizione strettamente positive ($T > 0$), si ha sempre $|Q_1| > 0$. Pertanto \textbf{è impossibile} che in un ciclo termico il calore ceduto alla sorgente fredda sia nullo. Questo fornisce un fondamento matematico rigoroso all'enunciato di Kelvin-Planck.
\end{osservazione}
\section{Entropia e Disuguaglianza di Clausius}

La \textbf{disuguaglianza di Clausius} fornisce una formulazione quantitativa del Secondo Principio, permettendo di distinguere analiticamente i cicli reversibili da quelli irreversibili.
Partendo da una macchina termica operante tra due sorgenti, si ha:
\begin{align}
    \text{Reversibile:} &\quad \frac{Q_1}{T_1} + \frac{Q_2}{T_2} = 0 \quad (\eta = \eta_C) \\
    \text{Irreversibile:} &\quad \frac{Q_1}{T_1} + \frac{Q_2}{T_2} < 0 \quad (\eta < \eta_C)
\end{align}
Generalizzando a una macchina che scambia calore con $n$ sorgenti termiche (o integrando per trasformazioni continue), si ottiene il teorema di Clausius:
\begin{equation}
    \boxed{\oint \frac{\delta Q}{T} \le 0}
\end{equation}
dove l'uguaglianza vale se e solo se l'intero ciclo è reversibile.

\vspace{0.3cm}
\subsection{Definizione di Entropia}
Se consideriamo una trasformazione \textbf{ciclica e reversibile}, l'integrale nullo implica che:
\begin{equation*}
    \oint_{\text{rev}} \frac{\delta Q}{T} = \int_{A}^{B} \left(\frac{\delta Q}{T}\right)_{\text{rev}} + \int_{B}^{A} \left(\frac{\delta Q}{T}\right)_{\text{rev}} = 0 \implies \int_{A}^{B} \left(\frac{\delta Q}{T}\right)_{\text{rev}} = \int_{A}^{B} \left(\frac{\delta Q}{T}\right)_{\text{rev, cammino 2}}
\end{equation*}
L'integrale lungo un cammino reversibile \textbf{non dipende dal percorso} ma solo dagli stati iniziale e finale. L'integrando $\frac{\delta Q}{T}$ risulta quindi un differenziale esatto. Questa proprietà definisce una nuova funzione di stato, chiamata \textbf{Entropia ($S$)}:
\begin{equation}
    \boxed{ \Delta S = S(B) - S(A) = \int_A^B \left( \frac{\delta Q}{T} \right)_{\text{rev}} }
\end{equation}
(L'unità di misura è J/K). 
Per un processo \textbf{irreversibile} da $A$ a $B$, si chiude il ciclo tornando ad $A$ con un processo reversibile ausiliario:
\begin{equation*}
    \oint \frac{\delta Q}{T} = \int_{A}^{B} \left( \frac{\delta Q}{T} \right)_{\text{irr}} + \int_{B}^{A} \left( \frac{\delta Q}{T} \right)_{\text{rev}} < 0
\end{equation*}
Da cui si ottiene: $\int_{A}^{B} \left( \frac{\delta Q}{T} \right)_{\text{irr}} < S(B) - S(A)$. In generale:
\begin{equation}
    \boxed{ \Delta S \ge \int_{A}^{B} \frac{\delta Q}{T} } \quad \text{(Disuguaglianza per trasformazioni aperte)}
\end{equation}

\begin{figure}[ht]
\centering
\begin{tikzpicture}
% Assi
\draw[->] (0,0) -- (3,0) node[right] {$V$};
\draw[->] (0,0) -- (0,3) node[above] {$P$};
\fill (0.5, 2) circle (2pt) node[above] {$A$};
\fill (2.5, 0.5) circle (2pt) node[right] {$B$};
\draw[->, dashed] (0.5, 2) to[out=-20, in=120] (1.5, 1.2) node[above] {I} to[out=-60, in=170] (2.5, 0.5);
\draw[<-, dashed] (0.5, 2) to[out=-80, in=160] (1.5, 0.8) node[below] {II} to[out=-20, in=180] (2.5, 0.5);
\end{tikzpicture}
\end{figure}

Per un processo \textbf{irreversibile} da $A$ a $B$, confrontando il percorso irreversibile con uno reversibile:
\begin{align*}
    \int_{A}^{B} \left( \frac{\delta Q}{T} \right)_{\text{irr}} &< S(B) - S(A) \\
    \int_{A}^{B} \frac{\delta Q}{T} &\le S(B) - S(A)
\end{align*}

\begin{definizione}{Formulazione Entropica del Secondo Principio della Termodinamica}
Per qualsiasi trasformazione da uno stato $A$ ad uno stato $B$, l'integrale di Clausius è sempre minore o uguale alla variazione di entropia del sistema. L'uguaglianza vale solo per trasformazioni reversibili.
\end{definizione}

\begin{definizione}{Variazione di Entropia per Sorgenti e Termostati}
Un termostato (o sorgente termica) è per definizione un serbatoio di calore a temperatura costante $T_0$ con capacità termica infinita ($C \to \infty$). Poiché lo scambio termico avviene idealmente in modo isotermo, senza gradienti interni di temperatura, la trasformazione per la sorgente si considera sempre reversibile. La sua variazione di entropia quando scambia un calore $Q_{\text{sorgente}}$ è calcolata elementarmente come:
\begin{equation}
    \boxed{ \Delta S_{\text{sorgente}} = \int \frac{\delta Q_{\text{sorgente}}}{T_0} = \frac{Q_{\text{sorgente}}}{T_0} }
\end{equation}
Se una sorgente cede una quantità di calore $Q > 0$ ad un sistema operativo, rispetto al termostato il calore è uscente ($Q_{\text{sorgente}} = -Q$), con conseguente diminuzione di entropia $\Delta S_{\text{sorgente}} = -\frac{Q}{T_0}$.
\end{definizione}

\vspace{0.3cm}
\subsection{Proprietà dell'Entropia}
\begin{itemize}
    \item \textbf{Additività}: un sistema composto da sottosistemi debolmente interagenti ha l'entropia pari alla somma delle singole entropie dei sottosistemi.
    \item \textbf{Estensività}: l'entropia totale di un sistema composto da più sotto-sistemi non interagenti a lungo raggio è la somma delle entropie delle singole parti. Essendo $\delta Q = \sum_i \delta Q_i$ e $T$ la medesima in condizioni di equilibrio termico, si ha:
    \begin{equation}
        \Delta S = \int_A^B \frac{\delta Q}{T} = \sum_i \int_A^B \frac{\delta Q_i}{T} = \sum_i \Delta S_i
    \end{equation}
    
    \item \textbf{Principio di Non Diminuzione dell'Entropia}: in un \textbf{sistema isolato}, non vi è alcuno scambio termico con l'esterno ($\delta Q = 0$). Pertanto:
    \begin{equation}
        \boxed{ \Delta S_{\text{isolato}} \ge 0 }
    \end{equation}
    
    \item \textbf{Entropia dell'Universo}: L'universo termodinamico, definito come l'insieme di Sistema e Ambiente, è per definizione un sistema isolato. Di conseguenza:
    \begin{equation*}
        \Delta S_{\text{univ}} = \Delta S_{\text{sist}} + \Delta S_{\text{amb}} \ge 0
    \end{equation*}
\end{itemize}

\begin{figure}[ht]
\centering
\begin{tikzpicture}
    \draw[->, >=stealth, thick, mypen] (0,0) -- (6,0) node[right] {Tempo $t$};
    \draw[->, >=stealth, thick, mypen] (0,0) -- (0,3.5) node[above] {Entropia $S$};
    \draw[dashed, thick, notered] (0, 3) node[left] {$S_{max}$} -- (5.8, 3);
    \draw[thick, blue] (0.2, 0.5) to[out=60, in=185] (1.5, 2.0) to[out=5, in=190] (3.0, 2.7) to[out=10, in=180] (5.5, 3.0);
    \node[align=center] at (3, 1) {\small I sistemi isolati evolvono spontaneamente\\ \small verso lo stato di massima entropia};
\end{tikzpicture}
\caption{Andamento nel tempo dell'entropia di un sistema isolato.}
\end{figure}

\subsection{Calcolo della Variazione di Entropia}
Poiché l'entropia è una funzione di stato, la sua variazione $\Delta S$ lungo una trasformazione \textbf{irreversibile} si calcola immaginando un cammino \textbf{reversibile arbitrario} che congiunge gli stessi stati $A$ e $B$.
\begin{itemize}
    \item \textbf{Espansione libera di un gas perfetto}: $\Delta S_{\text{gas}} = n R \ln(V_B/V_A) > 0$.
    \item \textbf{Trasformazione Isoterma}: $\Delta S = Q_{AB}/T = n R \ln(V_B/V_A)$.
    \item \textbf{Trasformazione Adiabatica Reversibile}: $\Delta S = 0$ (trasformazione isoentropica).
    \item \textbf{Trasformazione generica di un gas perfetto}: 
    \begin{align}
        \Delta S &= n c_v \ln\left(\frac{T_B}{T_A}\right) + n R \ln\left(\frac{V_B}{V_A}\right) = n c_p \ln\left(\frac{T_B}{T_A}\right) - n R \ln\left(\frac{p_B}{p_A}\right) \\
        &= n c_v \ln\left( \frac{T_B V_B^{\gamma-1}}{T_A V_A^{\gamma-1}} \right) = n c_v \ln\left( \frac{p_B V_B^\gamma}{p_A V_A^\gamma} \right)
    \end{align}
    \item \textbf{Passaggi di stato (a $T$ e $p$ costanti)}: $\Delta S = \frac{m \lambda}{T}$.
    \item \textbf{Solidi e liquidi incomprimibili}: $\Delta S = m c \ln\left(\frac{T_B}{T_A}\right)$.
\end{itemize}

\vspace{0.3cm}
Inoltre, per un gas perfetto è possibile definire la funzione di stato \textbf{Entropia Assoluta} (a meno di una costante additiva di riferimento $S_0$):
\begin{equation}
    S_{\text{gas}}(p, V) = n c_v \ln\left(p V^\gamma\right) + \text{costante}
\end{equation}
Mentre per un solido o liquido a calore specifico $C$ costante: $S(T) = C \ln T + \text{costante}$.

\begin{osservazione}{Esempio: Calcolo dell'Entropia in una Espansione Libera ed Equivalenza del Percorso Ausiliario}
Consideriamo una mole ($n = 1\text{ mol}$) di gas perfetto che si espande liberamente nel vuoto raddoppiando il proprio volume ($V_B = 2V_A$) all'interno di un contenitore rigido adiabatico ($Q=0, L=0$). Come già dimostrato, la temperatura finale coincide con quella iniziale ($T_B = T_A = T_0$). Poiché la trasformazione reale è irreversibile, per calcolare la variazione di entropia del gas $\Delta S_{\text{gas}}$ dobbiamo connettere gli stati di equilibrio $A$ e $B$ tramite una trasformazione reversibile ausiliaria, ad esempio un'\textbf{isoterma reversibile}:
\begin{equation}
    \Delta S_{\text{gas}} = \int_A^B \left(\frac{\delta Q}{T}\right)_{\text{rev}} = \frac{1}{T_0} \int_{V_A}^{2V_A} p \, dV = n R \ln\left(\frac{2V_A}{V_A}\right) = R \ln 2 \simeq 5.76 \, \frac{\text{J}}{\text{K}} > 0
\end{equation}
Per l'ambiente, nella trasformazione reale non vi è stato alcuno scambio termico ($Q_{\text{amb}} = 0 \implies \Delta S_{\text{amb}} = 0$). Di conseguenza, per l'universo:
\begin{equation}
    \Delta S_{\text{univ}} = \Delta S_{\text{gas}} + \Delta S_{\text{amb}} = R \ln 2 + 0 > 0
\end{equation}
Se volessimo ripristinare il gas allo stato iniziale comprimendolo reversibilmente a contatto con un termostato ausiliario a temperatura $T_0$, il lavoro esterno compiuto sul gas verrebbe integralmente ceduto come calore alla sorgente ($Q_{\text{sorgente}} = +T_0 \Delta S_{\text{gas}}$), incrementando permanentemente l'entropia del termostato di $+R \ln 2$. Ciò conferma l'impossibilità di cancellare dall'universo l'irreversibilità generata dall'espansione libera.
\end{osservazione}

\vspace{0.3cm}
\subsection{Il Diagramma Temperatura-Entropia (T-S)}
Il diagramma \textbf{Temperatura-Entropia (T-S)} rappresenta un fondamentale strumento analitico in termodinamica. Poiché per una trasformazione reversibile vale la definizione infinitesima $\delta Q_{\text{rev}} = T\,dS$, l'area sottesa alla curva di una trasformazione in tale piano rappresenta geometricamente il \textbf{calore scambiato}:
\begin{equation}
    \boxed{ Q_{\text{rev}} = \int_A^B T\,dS }
\end{equation}
Per un ciclo termodinamico reversibile chiuso, essendo la variazione di energia interna nulla ($\oint dU = 0$), l'area interna alla linea chiusa del ciclo rappresenta esattamente il \textbf{lavoro netto} compiuto: $L = \oint \delta Q = \oint T\,dS$.

\vspace{0.3cm}
\subsubsection{Rappresentazione delle Trasformazioni e Confronto delle Pendenze}
Nel piano T-S, le quattro trasformazioni fondamentali assumono forme geometriche caratteristiche:
\begin{itemize}
    \item \textbf{Isoterma ($T = \text{cost}$)}: è rappresentata da una segmento rettilineo \textbf{orizzontale}.
    \item \textbf{Adiabatica Reversibile / Isoentropica ($S = \text{cost}$)}: è rappresentata da un segmento rettilineo \textbf{verticale} ($\Delta S = 0$).
    \item \textbf{Isocora e Isobara}: sono rappresentate da curve esponenziali crescenti della forma $T(S) = T_0 e^{\Delta S / n c}$.
\end{itemize}

\begin{figure}[ht]
\centering
\begin{minipage}{0.46\textwidth}
\centering
\begin{tikzpicture}[scale=1.05]
    % Griglia
    \draw[step=1cm, gray!15, very thin] (0,0) grid (4.2,4.2);
    % Assi
    \draw[->, >=stealth, thick, mypen] (0,0) -- (0,4.4) node[left] {$T$};
    \draw[->, >=stealth, thick, mypen] (0,0) -- (4.4,0) node[below] {$S$};
    
    % Punto iniziale comune A(1.2, 1.2)
    % Isoterma (T = 1.2)
    \draw[thick, blue] (1.2, 1.2) -- (4.0, 1.2) node[right] {\small Isoterma};
    
    % Adiabatica (S = 1.2)
    \draw[thick, notered] (1.2, 1.2) -- (1.2, 4.0) node[above] {\small Adiabatica};
    
    % Isocora: T = 1.2 * exp((S-1.2)/1.2) -> ripida
    \draw[thick, mygreen, domain=1.2:2.65, samples=30] plot (\x, {1.2*exp((\x-1.2)/1.2)}) node[above] {\small Isocora};
    
    % Isobara: T = 1.2 * exp((S-1.2)/2.0) -> meno ripida
    \draw[thick, orange, domain=1.2:3.6, samples=30] plot (\x, {1.2*exp((\x-1.2)/2.0)}) node[right] {\small Isobara};
    
    % Punto A
    \filldraw[mypen] (1.2, 1.2) circle (1.8pt) node[below left] {$A(S_0, T_0)$};
    \draw[dashed, gray] (1.2,0) node[below] {$S_0$} -- (1.2,1.2) -- (0,1.2) node[left] {$T_0$};
\end{tikzpicture}
\end{minipage}\hfill
\begin{minipage}{0.50\textwidth}
\textbf{Confronto analitico delle pendenze}:\\
Dalla relazione differenziale del calore scambiato $\delta Q = n c\,dT = T\,dS$, si ricava la pendenza delle curve nel diagramma T-S:
\begin{equation}
    \left( \frac{dT}{dS} \right) = \frac{T}{n c}
\end{equation}
Confrontando una trasformazione isocora ($c = c_v$) con una isobara ($c = c_p$) passanti per lo stesso stato $(S_0, T_0)$:
\begin{equation}
    \left( \frac{dT}{dS} \right)_v = \frac{T}{n c_v} > \frac{T}{n c_p} = \left( \frac{dT}{dS} \right)_p
\end{equation}
Poiché $c_p = c_v + R > c_v$, nel piano T-S \textbf{la curva isocora cresce sempre con una pendenza maggiore rispetto alla curva isobara}.
\end{minipage}
\caption{Confronto nel piano Temperatura-Entropia tra le quattro trasformazioni termodinamiche elementari a partire dal medesimo stato $A$.}
\end{figure}

\vspace{0.3cm}
\subsubsection{Il Ciclo di Carnot nel Piano T-S}
La potenza concettuale del diagramma T-S si manifesta pienamente nell'analisi del \textbf{Ciclo di Carnot}. Poiché è composto esclusivamente da due isoterme (a $T_2$ e $T_1$) e due adiabatiche reversibili (isoentropiche a $S_B$ e $S_A$), nel piano T-S il ciclo di Carnot assume la geometria elementare di un \textbf{rettangolo}.

\begin{figure}[ht]
\centering
\begin{minipage}{0.44\textwidth}
\centering
\begin{tikzpicture}[scale=1.1]
    % Griglia
    \draw[step=1cm, gray!15, very thin] (0,0) grid (4.2,3.8);
    % Assi
    \draw[->, >=stealth, thick, mypen] (0,0) -- (0,4.0) node[left] {$T$};
    \draw[->, >=stealth, thick, mypen] (0,0) -- (4.2,0) node[below] {$S$};
    
    % Rettangolo di Carnot tra S_A=1.2, S_B=3.4 e T1=1.0, T2=3.0
    \fill[cyan!15, opacity=0.8] (1.2, 1.0) rectangle (3.4, 3.0);
    
    % Ciclo
    \draw[thick, blue, ->, >=stealth] (1.2, 3.0) -- (2.4, 3.0);
    \draw[thick, blue] (2.4, 3.0) -- (3.4, 3.0) node[above, midway] {$Q_2 > 0$};
    
    \draw[thick, notered, ->, >=stealth] (3.4, 3.0) -- (3.4, 2.0);
    \draw[thick, notered] (3.4, 2.0) -- (3.4, 1.0);
    
    \draw[thick, blue, ->, >=stealth] (3.4, 1.0) -- (2.3, 1.0);
    \draw[thick, blue] (2.3, 1.0) -- (1.2, 1.0) node[below, midway] {$Q_1 < 0$};
    
    \draw[thick, notered, ->, >=stealth] (1.2, 1.0) -- (1.2, 2.0);
    \draw[thick, notered] (1.2, 2.0) -- (1.2, 3.0);
    
    % Vertici
    \filldraw[mypen] (1.2, 3.0) circle (1.6pt) node[above left] {$A$};
    \filldraw[mypen] (3.4, 3.0) circle (1.6pt) node[above right] {$B$};
    \filldraw[mypen] (3.4, 1.0) circle (1.6pt) node[below right] {$C$};
    \filldraw[mypen] (1.2, 1.0) circle (1.6pt) node[below left] {$D$};
    
    % Coordinate
    \draw[dashed, gray] (1.2, 0) node[below] {$S_A$} -- (1.2, 1.0);
    \draw[dashed, gray] (3.4, 0) node[below] {$S_B$} -- (3.4, 1.0);
    \draw[dashed, gray] (0, 3.0) node[left] {$T_2$} -- (1.2, 3.0);
    \draw[dashed, gray] (0, 1.0) node[left] {$T_1$} -- (1.2, 1.0);
    
    \node[mypen, font=\bfseries] at (2.3, 2.0) {$L > 0$};
\end{tikzpicture}
\end{minipage}\hfill
\begin{minipage}{0.52\textwidth}
Dal grafico rettangolare si deducono per immediata via geometrica tutte le grandezze del ciclo:
\begin{itemize}
    \item \textbf{Calore assorbito} ($A \to B$ a $T_2$ costante):
    \begin{equation}
        Q_2 = \int_{S_A}^{S_B} T_2\,dS = T_2 (S_B - S_A) > 0
    \end{equation}
    \item \textbf{Calore ceduto} ($C \to D$ a $T_1$ costante):
    \begin{equation}
        Q_1 = \int_{S_B}^{S_A} T_1\,dS = - T_1 (S_B - S_A) < 0
    \end{equation}
    \item \textbf{Lavoro netto} (Area del rettangolo azzurro):
    \begin{equation}
        L = Q_2 - |Q_1| = (T_2 - T_1)(S_B - S_A)
    \end{equation}
\end{itemize}
Il rendimento del ciclo si ottiene quindi con una banale semplificazione algebrica:
\begin{equation}
    \eta_C = \frac{L}{Q_2} = \frac{(T_2 - T_1)(S_B - S_A)}{T_2 (S_B - S_A)} = 1 - \frac{T_1}{T_2}
\end{equation}
\end{minipage}
\caption{Il ciclo di Carnot rappresentato nel piano T-S: le due isoterme orizzontali e le due adiabatiche verticali delimitano un perfetto rettangolo, rendendo immediata la deduzione geometrica del rendimento.}
\end{figure}

\subsection{Degradazione dell'Energia e Lavoro Perduto}
Consideriamo due cicli operanti fra gli stessi termostati $T_A, T_B$, di cui uno di Carnot ($M_C$) e uno irreversibile ($M$).

\begin{figure}[ht]
\centering
\begin{tikzpicture}[scale=1.1]
% Sorgente Calda
\fill[pattern=north east lines, notered!25] (0, 3) rectangle (3, 4);
\draw[thick, mypen, fill=orange!15] (0, 3) rectangle (3, 4) node[midway, font=\bfseries] {Sorgente Calda ($T_B$)};
\node[left, font=\bfseries, notered] at (0, 3.5) {$T_B > T_A$};

% Macchina Termica
\draw[thick, mypen, fill=cyan!15] (1.5, 1.5) circle (0.6);
\node[font=\bfseries] at (1.5, 1.5) {$M / M_C$};

% Sorgente Fredda
\fill[pattern=north west lines, cyan!25] (0, -1) rectangle (3, 0);
\draw[thick, mypen, fill=cyan!10] (0, -1) rectangle (3, 0) node[midway, font=\bfseries] {Sorgente Fredda ($T_A$)};
\node[left, font=\bfseries, blue] at (0, -0.5) {$T_B > T_A$};

% Frecce di scambio termico
\draw[->, >=stealth, thick, notered, line width=1.5pt] (1.5, 3) -- (1.5, 2.1) node[midway, right, font=\bfseries] {$Q_B$};
\draw[->, >=stealth, thick, blue, line width=1.5pt] (1.5, 0.9) -- (1.5, 0) node[midway, right, font=\bfseries] {$Q_A$ (o $Q_{AC}$)};

% Lavoro W
\draw[->, >=stealth, thick, mygreen, line width=1.5pt] (2.1, 1.5) -- (3.5, 1.5) node[right, font=\bfseries] {$L_M$ (o $L_{M_C}$)};
\end{tikzpicture}
\caption{Confronto termodinamico tra un motore reversibile di Carnot ($M_C$) e un motore reale irreversibile ($M$) operanti tra i medesimi serbatoi termici $T_B > T_A$: a parità di calore assorbito $Q_B$, il motore irreversibile cede maggior calore alla sorgente fredda ($|Q_A| > |Q_{AC}|$), comportando una perdita di lavoro utile ($L_M < L_{M_C}$) proporzionale all'entropia generata.}
\end{figure}

Entrambe le macchine $M$ e $M_C$ assorbono lo stesso calore $Q_B$ dalla sorgente calda.
Tuttavia, esse cedono calore alla sorgente fredda in modo diverso: $M$ cede $Q_A$, mentre $M_C$ cede $Q_{AC}$.
Essendo $M_C$ reversibile, il suo rendimento è quello di Carnot:
\begin{equation*}
    \eta_C = 1 - \frac{|Q_{AC}|}{Q_B} = 1 - \frac{T_A}{T_B}
\end{equation*}
Per la macchina irreversibile $M$, il rendimento è strettamente minore di quello di Carnot:
\begin{equation*}
    \eta = 1 - \frac{|Q_A|}{Q_B} < 1 - \frac{T_A}{T_B}
\end{equation*}

Poiché entrambe compiono cicli completi, la variazione di entropia delle macchine è nulla ($\Delta S_M = \Delta S_{M_C} = 0$). 
La variazione di entropia dell'Universo per il ciclo irreversibile è quindi data solo dalle sorgenti:
\begin{equation*}
    \Delta S_{\text{univ}} = \Delta S_{T_A} + \Delta S_{T_B} = - \frac{Q_A}{T_A} - \frac{Q_B}{T_B} > 0
\end{equation*}
Mentre per il ciclo reversibile $M_C$ risulta $\Delta S_{\text{univ}} = 0$.

Possiamo ora calcolare la differenza tra il lavoro prodotto dalla macchina reversibile ($L_{M_C}$) e quello prodotto dalla macchina reale ($L_M$):
\begin{align*}
    L_M &= Q_A + Q_B = (Q_A - Q_{AC}) + (Q_{AC} + Q_B) \\
    &= (Q_A - Q_{AC}) + L_{M_C}
\end{align*}
Sostituendo $Q_{AC} = - \frac{T_A}{T_B} Q_B$ (dal ciclo di Carnot), otteniamo:
\begin{align*}
    L_M &= L_{M_C} + Q_A + \frac{T_A}{T_B} Q_B = L_{M_C} + T_A \left( \frac{Q_A}{T_A} + \frac{Q_B}{T_B} \right) \\
    &= L_{M_C} - T_A \Delta S_{\text{univ}}
\end{align*}
Da cui si ricava infine che:
\begin{equation}
    \boxed{ L_M = L_{M_C} - T_A \Delta S_{\text{univ}} < L_{M_C} }
\end{equation}
La differenza $\boxed{ L_{M_C} - L_M = T_A \Delta S_{\text{univ}} }$ esprime il teorema di \textbf{Gouy-Stodola}: il lavoro meccanico perduto a causa delle irreversibilità è direttamente proporzionale all'entropia generata nell'universo.

\subsection{Casi Tipici di Irreversibilità}
L'irreversibilità di un processo in un sistema isolato si dimostra verificando che la variazione di entropia dell'Universo (Sistema + Ambiente) sia strettamente positiva ($\Delta S_{\text{univ}} > 0$).

\begin{itemize}
    \item \textbf{CASO 1: Riscaldamento per attrito} \\
    Un corpo è trascinato su una superficie scabra. Tutto il lavoro speso per vincere l'attrito si trasforma in variazione di energia interna ($\Delta U = L_{\text{attrito}}$), innalzando la temperatura del sistema (corpo + superficie). Essendo isolato dall'esterno ($Q = 0$), l'ambiente non subisce variazioni. Per calcolare l'entropia, si congiungono gli stati mediante un riscaldamento isocoro reversibile:
    \begin{equation}
        \Delta S_{\text{univ}} = \Delta S_{\text{sist}} = C \ln\left( \frac{T_f}{T_i} \right) > 0 \quad (\text{poiché } T_f > T_i)
    \end{equation}

    \item \textbf{CASO 2: Espansione libera di un gas} \\
    Come analizzato in precedenza, l'espansione nel vuoto di un gas perfetto ($Q=0, L=0, \Delta U=0$) avviene a temperatura costante, ma il volume aumenta. Poiché il sistema è isolato termicamente:
    \begin{equation}
        \Delta S_{\text{univ}} = \Delta S_{\text{sist}} = nR \ln\left( \frac{V_B}{V_A} \right) > 0 \quad (\text{poiché } V_B > V_A)
    \end{equation}

    \item \textbf{CASO 3: Scambio termico per conduzione} \\
    Consideriamo due corpi solidi isolati, con uguale capacità termica $C$ ma temperature iniziali $T_1 \neq T_2$. Essi vengono posti in contatto e raggiungono l'equilibrio a $T_f = \frac{T_1 + T_2}{2}$. L'entropia dell'universo (che coincide col sistema dei due corpi) è la somma delle due variazioni, calcolate su isocore reversibili:
    \begin{equation}
        \Delta S_{\text{univ}} = C \ln\left( \frac{T_f}{T_1} \right) + C \ln\left( \frac{T_f}{T_2} \right) = C \ln\left( \frac{T_f^2}{T_1 T_2} \right)
    \end{equation}
    Ricordando che la media aritmetica è strettamente maggiore della media geometrica per numeri disuguali, si ha $T_f^2 > T_1 T_2$. Dunque il logaritmo è positivo e $\Delta S_{\text{univ}} > 0$.

    \item \textbf{CASO 4: Ciclo monotermo} \\
    In un ciclo termodinamico che scambia calore con una sola sorgente a temperatura $T$, il sistema torna allo stato iniziale e dunque $\Delta S_{\text{sist}} = 0$. Il Primo Principio impone che il lavoro eseguito sia pari al calore assorbito ($L = Q$). Per il principio di Kelvin-Planck, tale lavoro non può essere positivo ($L \le 0 \implies Q \le 0$). Il calore assorbito dalla sorgente è quindi $-Q \ge 0$, da cui:
    \begin{equation}
        \Delta S_{\text{univ}} = \Delta S_{\text{sist}} + \Delta S_{\text{amb}} = 0 - \frac{Q}{T} \ge 0
    \end{equation}
    Se il ciclo produce attriti e richiede lavoro per essere completato ($L < 0$), allora $\Delta S_{\text{univ}} > 0$ e il ciclo è irreversibile.
\end{itemize}

\vspace{0.5cm}
\begin{flushright}
\textit{Fine del Capitolo 6 -- Termodinamica}
\end{flushright}
